\documentclass[11pt]{article}
\usepackage[margin=1in]{geometry}
\usepackage{amsmath,amssymb,amsthm,mathtools,booktabs,longtable,array}
\usepackage[T1]{fontenc}
\usepackage{lmodern,microtype}
\usepackage[colorlinks=true,allcolors=blue,bookmarksnumbered=true]{hyperref}
\numberwithin{equation}{section}
\newtheorem{theorem}{Theorem}[section]
\newtheorem{lemma}[theorem]{Lemma}
\newtheorem{proposition}[theorem]{Proposition}
\newtheorem{corollary}[theorem]{Corollary}
\theoremstyle{definition}
\newtheorem{definition}[theorem]{Definition}
\newtheorem{assumption}[theorem]{Assumption}
\newtheorem{example}[theorem]{Example}
\theoremstyle{remark}
\newtheorem{remark}[theorem]{Remark}
\newcommand{\tr}{\operatorname{tr}}
\newcommand{\id}{\operatorname{id}}
\newcommand{\dd}{\,\mathrm d}
\setlength{\emergencystretch}{2em}
\title{Complete-Current Estimates for Coherent\\Quantum Sources and Continua}
\author{Jeremy Rodgers\thanks{Website: \href{https://everythingequation.com}{everythingequation.com}}\\Independent Researcher}
\date{9 October 2026\\[3pt]
\small\href{https://doi.org/10.5281/zenodo.23259571}{doi:10.5281/zenodo.23259571}}
\hypersetup{pdftitle={Complete-Current Estimates for Coherent Quantum Sources and Continua},
pdfauthor={Jeremy Rodgers},pdfsubject={Coherent response, complete currents and retained quantum sources}}
\begin{document}
\maketitle
\begin{abstract}
We prove quantitative response estimates that retain coherent quantum
sources, complete physical currents, spectral poles and actual feedback.
An ordered inverse lift converts temporally regular form-dual forcing
into a first-form response under common-form work bounds.
A finite spectral-window theorem controls varying Hilbert-space inputs,
and ordered closed-form resolvent identities transfer noncommuting
physical-energy weights. Far-frequency and damped estimates retain the
endpoint and source-graph costs. Complementary force, invariant and
temporal-domain estimates distinguish valid current resources for
singular interactions. In a compact-profile electron--oscillator model,
we derive the full Coulomb forcing measure and an actual-source
low-window sandwich below $4.05\times10^{-9}$. A separate finite-time
calculation bounds the electron and source current differences by
$.046$ and $.004$ in their respective coordinate lengths.
The spectral supremum is taken over the actual spectrum; a spectral
floor supplies only an upper bound. Full-source forcing compression
and causal feedback are developed in appendices.
These results provide explicit analytic inputs to physical transport
and record problems. Event probabilities and a realized measurement
apparatus require additional flow, preparation and decoder premises.
\end{abstract}
\clearpage
\tableofcontents
\newpage
\section{Introduction}
\label{p4s:sec:introduction}

An error estimate for a quantum amplitude becomes a physical transport
estimate only after the relevant derivatives, source coordinates and
current constitution have been retained. This distinction is especially
important for a coherently driven continuum. A small omitted occupation
does not eliminate its interference current; a small omitted forcing
does not eliminate response generated by feedback; and an accurately
resolved electronic transition does not determine the current of a
quantum source. These are analytic questions even before an experiment
or a measurement record is specified.

We develop a common set of quantitative interfaces for these questions.
The complete wave remains a vector on all admitted positions and
internal degrees of freedom. Its error is controlled in a form that
dominates the physical kinetic derivatives, then converted to the full
coherent current. The response may be driven by an unbounded coupling,
contain genuine bound poles and a continuum, and receive a
time-dependent vector input from another interacting block. We retain
the input of the coupled evolution throughout.

There are three main results. First, an ordered inverse lift converts
temporally regular form-dual forcing into a first-form response under
absolutely continuous common-form work bounds. It avoids a weighted
inverse-square-root derivative that is not uniformly controlled in
dimension. Second, an operator spectral-window estimate controls a
varying Hilbert-space input. The proof uses interval means, spectral
derivatives and orthogonal output sectors; it does not replace the
vector input by a scalar supremum. Ordered resolvent identities then
transfer genuinely noncommuting physical-energy weights, while explicit
far-frequency dressings retain both endpoints. Third, a compact-profile
electron--oscillator model supplies an analytic Coulomb spectral
calculation, an actual-source inverse bound and a complete finite-time
current calculation with source feedback.

Several complementary results make these interfaces usable. A spatial
force estimate retains covariant strain, curvature and first-order
transport. A temporal common-domain estimate instead accepts a residual
without a spatial form norm and retains its right clock. An exact
positive quadratic invariant supplies a comparison form even through
an inverted trap interval. The appendices preserve a spatial forcing
compression with arbitrary source coefficients and the causal
high-sector feedback that such compression must not erase. Each has
its own input hypotheses; choosing the smallest coefficient from
incompatible estimates would not be a valid composition.

\subsection{Relation to established methods and prior work}

The functional-analytic ingredients have substantial precedents.
Kato's smoothing and factorized-resolvent framework
\cite{Kato1966} and its inhomogeneous development
\cite{DAncona2014} motivate the relation between response and resolvents.
Our fixed positive imaginary part is an earned finite-resolution
coefficient, not a claim of a uniform limiting-absorption estimate.
Common-form propagation belongs to the classical evolution theory of
Kisy\'nski \cite{Kisynski1964}; we give the argument for the precise
absolutely continuous relative-work assumptions used here.
Common \emph{operator}-domain regularity, used in the separate temporal
graph theorem, is discussed in \cite{SchmidGriesemer2014}.
We keep these two kinds of domains distinct.

The oscillator invariant is classical
\cite{Lewis1967,LewisRiesenfeld1969}; its contribution here is an explicit
physical-current metric with the centroid and mass factors preserved.
The Coulomb functions and their normalization are standard
\cite{DLMFCoulomb}; the spectral calculation is included to make the
forcing measure and source-weighted bounds reproducible. Isolated-band
constructions such as \cite{PanatiSpohnTeufel2003} require their own
symbol classes and gap hypotheses. Neither a finite column corrector
nor a frozen internal gap is asserted here to supply such a band
reduction of the complete position/source parent.

The current October research editions of the measurement programme
include equilibrium and nonequilibrium record constructions
\cite{RodgersEquilibrium,RodgersNonequilibrium}. Those works ask which
original laws yield outcome weights and whether a distinct receiver
preserves an earlier actual event. Their apparatus constants and
preparation premises are not inputs to the estimates proved here.
The new companion manuscripts treat retained-archive preparation
(P1)~\cite{RodgersP1}, effective repeated records (P2)~\cite{RodgersP2},
and reference-compatible flows and whole-history stability
(P3)~\cite{RodgersP3}. The present response and current estimates enter
such arguments only when their model, entrance and flow premises
agree. The proofs below are independent of those companions; the
specific flow application in Section~\ref{p4h:sec:flowinterface} uses
the revised P3 theorem explicitly.

These new manuscripts remain distinct from both the current October
website editions and the preserved September editions. No integration
into one apparatus is asserted. In particular, assumptions of the
pilot-medium and massive-configuration constitutions are not
interchangeable. The contribution here is the proved combination of
response, derivative and complete-current estimates under stated
hypotheses, rather than a priority claim for their established
individual ingredients.

\subsection{Conventions and the structure of the results}

We use $i\partial_t\Psi=H\Psi$ with $H$ in rate units, unless reduced-mass
atomic units are explicitly declared. Hilbert-space norms include every
physical source coordinate; a finite passive internal reference can be
included by tensoring the identity. A source wave or its quantum
covariance is not an assumption about an actual configuration law.
Form inequalities always refer to the complete stated parent.

Section~\ref{p4s:sec:current} first supplies the current and coordinate
conversion. Sections~\ref{p4f:sec:lift}--\ref{p4f:sec:derivatives}
prove the form-dual, spectral-window and ordered-resolvent results.
Section~\ref{p4s:sec:forms} develops the alternative spatial-force,
invariant and temporal-domain estimates.
Sections~\ref{p4h:sec:spectral}--\ref{p4h:sec:currents} give the
hydrogenic example. Appendices~\ref{p4s:sec:spatial} and
\ref{p4s:sec:feedback} preserve the full-source spatial compression
and feedback mechanisms.

The hydrogenic example has
$\Omega=.01$, $|g|\leq10^{-6}$ and a radius-$40$ compact dipole
profile in its specified units. Its low-energy physical resolvent
sandwich is below $4.05\times10^{-9}$ under the stated source weight.
A separate all-frequency, stronger-source-weight estimate gives, for
the declared entrance and $T=100$, integrated electron and source
\emph{current differences} below $.046$ and $.004$ in their respective
coordinate lengths. The comparison uses the actual projected wave,
including feedback. It is not a second independently flowing wave,
and these numbers are not probabilities.

\section{From a coherent error to the complete physical current}
\label{p4s:sec:current}

We first state the conversion that every response estimate in this paper
must supply. All positions, including source positions, belong to
$x\in\mathbb R^d$. A finite internal fibre is allowed. A wave can also be a
matrix of finitely many input columns; all norms below are the spatial
$L^2$ Hilbert--Schmidt norms in that case. This convention gives one
envelope before contraction with an unknown input, rather than a sum of
incoherent branch currents. Write
$\operatorname{Im}_{H} Z=(Z-Z^\dagger)/(2i)$.

In rate units, consider a self-adjoint local parent with
\begin{equation}\label{p4s:eq:parent}
 H=-\sum_i D_i\lambda_i D_i+U+V
       -i\sum_i\left(a_iD_i+\tfrac12\nabla_i a_i\right),
 \quad D_i=\partial_i-iC_i,\quad
 \nabla_i X=\partial_iX-i[C_i,X].
\end{equation}
Here $C_i,a_i,V$ are Hermitian matrices, $U$ is real scalar, and
$\lambda_i>0$ is independent of position; dependence on time is allowed.
Domains and coefficient regularity sufficient for the indicated
continuity equation are part of this parent. Its density and current are
\begin{equation}\label{p4s:eq:literalcurrent}
 \rho_\Psi=\Psi^\dagger\Psi,\qquad
 J_i[\Psi]=2\lambda_i\operatorname{Im}_{H}(\Psi^\dagger D_i\Psi)
                          +\Psi^\dagger a_i\Psi .
\end{equation}
The complete continuity equation is
$\partial_t\rho_\Psi+\sum_i\partial_iJ_i[\Psi]=0$; it also holds
entrywise for the matrix of input columns evolved by the same parent.
Indeed, on a common local core,
$\partial_t(\Psi^\dagger\Psi)=i(H\Psi)^\dagger\Psi-i\Psi^\dagger H\Psi$.
The Hermitian multiplication terms cancel. The covariant product rule
turns the kinetic contribution into
$-\sum_i\partial_i[2\lambda_i\operatorname{Im}_H(\Psi^\dagger D_i\Psi)]$
and the symmetrized first-order contribution into
$-\sum_i\partial_i(\Psi^\dagger a_i\Psi)$.
Testing and passage in the admitted form domain give the weak identity.
Thus the first-order term contributes a literal transport current;
omitting it would change the conservation law. Additional microscopic spin-curl or nonlocal
current constitutions require their own terms; they are not consequences
of \eqref{p4s:eq:parent}.

\begin{proposition}[Complete coherent current conversion]
\label{p4s:prop:current}
Let $\Psi=P+E$ use the same derivative and current constitution, and let
a positive form $K$ satisfy
$K\geq\sum_i\lambda_iD_i^\dagger D_i$. Put
\[
 e=\|E\|,\qquad Q=\|K^{1/2}E\|,\qquad
 N=\|\Psi\|,\qquad \tau_i=\|\sqrt{\lambda_i}D_iP\|.
\]
Whenever the displayed moments are finite,
\begin{align}
 \int\|\rho_\Psi-\rho_P\|_{\rm op}
   &\leq e(\|P\|+N),\label{p4s:eq:density}\\
 \int\|J_i[\Psi]-J_i[P]\|_{\rm op}
   &\leq2\sqrt{\lambda_i}(e\tau_i+NQ)
          +e\bigl(\|a_iP\|+\|a_i\Psi\|\bigr).
          \label{p4s:eq:currentbound}
\end{align}
If $a_i$ is bounded, the last term is at most
$\|a_i\|_\infty e(\|P\|+N)$. Neither normalization nor an
independent flow for the comparison $P$ is required.
\end{proposition}
\begin{proof}
Before any spatial integration, the exact identities are
\[
 \rho_\Psi-\rho_P=E^\dagger P+\Psi^\dagger E
\]
and
\begin{align*}
 J_i[\Psi]-J_i[P]
  ={}&2\lambda_i\operatorname{Im}_{H}
       \left(E^\dagger D_iP+\Psi^\dagger D_iE\right)\\
    &+E^\dagger a_iP+(a_i\Psi)^\dagger E .
\end{align*}
The operator norm of a product is bounded by the product of its
Hilbert--Schmidt norms, and
$\|\operatorname{Im}_{H}Z\|_{\rm op}\leq\|Z\|_{\rm op}$.
Spatial Cauchy--Schwarz and
$\sqrt{\lambda_i}\|D_iE\|\leq Q$ give the claims.
Hermiticity moves the unbounded $a_i$ in the last term onto $\Psi$,
so no unproved bounded-operator coefficient or moment of $a_iE$
has been inserted.
\end{proof}

For a charged particle with physical momentum
$p_{\rm ph}=-i\hbar\partial$, dividing its Hamiltonian by $\hbar$
gives $\lambda_i=\hbar/(2M_i)$ and
$C_i=e_iA_i/\hbar$. A symmetrized mechanical term
$\{p_{\rm ph},d\times B\}/(2M\hbar)$ in the rate generator has
$a=(d\times B)/M$: there is no additional factor $1/\hbar$ in its
velocity. An unbounded dipole or source operator must be handled by the
moments in \eqref{p4s:eq:currentbound}.

\begin{example}[Small wave error with unbounded current error]
\label{p4s:ex:norm}
On the circle of length $2\pi$, let
$\Psi=(2\pi)^{-1/2}$ and
\[
 P_n(t,x)=
 \frac{1+n^{-1}e^{in^2x-i\lambda n^4t}}
      {\sqrt{2\pi(1+n^{-2})}},\qquad n\geq2 .
\]
Both are normalized solutions of $i\partial_tu=-\lambda\partial_x^2u$,
and $\|P_n-\Psi\|\leq2/n$. Nevertheless
\[
 J[P_n]=\frac{2\lambda}{2\pi(1+n^{-2})}
              \left(n\cos(n^2x-\lambda n^4t)+1\right),
\]
so
$\|J[P_n]-J[\Psi]\|_{L^1}
\geq2\lambda(2n/\pi-1)/(1+n^{-2})\to\infty$.
The missing resource is derivative control, not a sharper wave-norm
inequality. This example even uses two legitimate free waves.
\end{example}

\subsection{Coordinates, moving tests and the original law}

\begin{proposition}[Current under a moving coordinate chart]
\label{p4s:prop:chart}
Let $(t,x)\mapsto\chi_t(x)$ be jointly $C^1$, with each
$q=\chi_t(x)$ an orientation-preserving diffeomorphism. Put
$G=D_x\chi_t$ and $g=\det G>0$. A locally integrable laboratory
density/current pair $(\rho,j)$ satisfying the continuity equation
transforms to
\begin{equation}\label{p4s:eq:chart}
 \widetilde\rho=g\,\rho\circ\chi_t,\qquad
 \widetilde j=gG^{-1}
       \left(j\circ\chi_t-(\rho\circ\chi_t)\partial_t\chi_t\right).
\end{equation}
These obey the transformed continuity equation. In particular a moving
surface $F(t,q)=n(t)\cdot(q-C(t))-s$, $|n|=1$, has relative normal flux
\begin{equation}\label{p4s:eq:surface}
 j\cdot n+\rho\,n'\cdot(q-C)-\rho\,n\cdot C'.
\end{equation}
\end{proposition}
\begin{proof}
For a compactly supported smooth test $\varphi(t,x)$, set
$\psi(t,q)=\varphi(t,\chi_t^{-1}(q))$. The chain rule gives
$\nabla_q\psi=G^{-T}\nabla_x\varphi$ and
$\psi_t=\varphi_t-\nabla_x\varphi\cdot G^{-1}\chi_t'$.
Insert these in the weak continuity equation and use $dq=g\,dx$.
This proves \eqref{p4s:eq:chart} without a componentwise guess at the
current. The derivative along a laboratory path is
$dF/dt=F_t+(j/\rho)\cdot\nabla F$; since $|\nabla F|=1$,
its density-weighted normal rate is \eqref{p4s:eq:surface}.
\end{proof}

The chart formula fixes the transport part of a transformed current. A
position-dependent internal unitary also transforms the connection;
one must conjugate the whole parent and include that derivative.
Integrating hidden source positions first is generally insufficient:
$|\int j\,dy|\leq\int|j|\,dy$ may discard counterflow or coherent
terms needed by a trajectory test.

To state the event interface precisely, suppose the complete current
already admits a reference path law with one-time density $\rho$ and
velocity $j/\rho$, and the original actual law satisfies
$\mu_0\leq C\rho_0$. Reweight that same path law by its entrance density
ratio. Admit tests $h$ for which $t\mapsto h(t,X_t)$ is absolutely
continuous on almost every reference path, with derivative
$h_t+(j/\rho)\cdot\nabla h$, and for which the integral below is finite.
For example, $C^1$ tests with this integrability have the required
chain rule. Then
\begin{equation}\label{p4s:eq:path}
 \mathbb E_{\mu_0}\operatorname{Var}_{[0,T]}h(t,X_t)
 \leq C\int_0^T\!\!\int
          |\rho\,h_t+j\cdot\nabla h|\,dq\,dt .
\end{equation}
Indeed, the chain rule along reference paths, entrance domination and
Tonelli's theorem prove this in that order. If
$h_t+b_{\rm ref}\cdot\nabla h=0$, the integrand becomes
$|\nabla h\cdot(j-b_{\rm ref}\rho)|$. Comparison to
$j_{\rm ref}=b_{\rm ref}\rho_{\rm ref}$ therefore costs both
$j-j_{\rm ref}$ and $b_{\rm ref}(\rho_{\rm ref}-\rho)$.
An event estimate additionally needs the test's guard width, initial bad
mass, interval, decoder and own reference crossings. Flow existence and
reference-compatible selection are separate premises, studied in
the companion flow paper~\cite{RodgersP3}. Equations
\eqref{p4s:eq:currentbound} and \eqref{p4s:eq:path} do not assign a new
actual law at an intermediate time. The domination premise concerns
the complete original joint law, including retained source and archive
coordinates. Absolute continuity alone supplies no finite value of $C$,
and an averaged conditional bound supplies no uniform guarantee after
arbitrarily rare postselection.

A sharp surface requires an appropriate trace and crossing area formula;
a volume $L^1$ current estimate or an almost-every-level coarea statement
does not by itself control flux at a prescribed surface. Smooth guard
tests can instead bound a traversal that changes $h$ by a stated positive
amount. Tangential contacts and moving guards must satisfy their own
chain-rule or crossing premises. All such estimates concern the full
time interval in \eqref{p4s:eq:path}.


\section{Temporal dressing in the physical form dual}
\label{p4f:sec:lift}

Response estimates must keep the input that actually drives the
complement. For an admitted fixed block decomposition
$\mathcal H_P\oplus\mathcal H_Q$, the equation
\[
 i\dot q=Aq+Lp
\]
uses the actual $p=P\Psi$, including its feedback through $q$,
source phases and clock dynamics. An incident population or a
frozen electronic amplitude is not a substitute for this vector.
Throughout this section energies are in the units of
$i\dot q=Aq+f$; a physical Hamiltonian must be divided by
$\hbar$ if time is measured in seconds.

Let $\mathcal V\subset\mathcal H\subset\mathcal V^*$ be a
dense separable Hilbert form triple, with the usual pivot
identification. A positive form $K\ge I$ identifies $\mathcal V$
with $\operatorname{Dom}K^{1/2}$ and gives the dual norm
$\|f\|_{\mathcal V^*,K}=\|K^{-1/2}f\|_{\mathcal H}$.
Here the square root on a dual vector is its continuous
extension, rather than an assertion that $f\in\mathcal H$.
We use the same letter for a self-adjoint operator and its
continuous form map $\mathcal V\to\mathcal V^*$ when the
domain makes the meaning unambiguous.

\begin{assumption}[Common physical form and relative work]
\label{p4f:ass:forms}
On $[0,T]$, $A(t)$ is associated with a closed Hermitian
form $a_t$ on the same $\mathcal V$. There is a real bounded
absolutely continuous scalar $c(t)$ such that
\[
 k_t=a_t+c(t)\langle\cdot,\cdot\rangle,\qquad K(t)=A(t)+c(t)\ge I.
\]
The $k_t$ norms are uniformly equivalent to one fixed form norm.
For fixed $u,v\in\mathcal V$, $k_t(u,v)$ is absolutely
continuous, with a strongly measurable form derivative and
its integral identity, satisfying
\begin{equation}\label{p4f:eq:relativework}
 |\dot k_t(u,v)|
 \le \nu(t)\|K(t)^{1/2}u\|\|K(t)^{1/2}v\|,
 \qquad \nu\ge0,\quad \nu\in L^1(0,T).
\end{equation}
Strong measurability here includes that the derivative applied
to any fixed form vector is measurable in $\mathcal V^*$.
\end{assumption}

The shift is an estimate of positivity, not a spectral gap at
an incident energy. Nor does this assumption require an
operator-norm derivative of a point-Coulomb force.
Classical common-form evolution under smoother form
hypotheses goes back to Kisy\'nski \cite{Kisynski1964}.
We give the conforming argument for the precise
absolutely continuous relative-work assumptions used here.

\begin{lemma}[Homogeneous common-form propagation]
\label{p4f:lem:propagation}
Under Assumption~\ref{p4f:ass:forms} there is a unique unitary
propagator $U(t,s)$ on $\mathcal H$. It preserves $\mathcal V$
and, for $s\le t$,
\begin{equation}\label{p4f:eq:propagation}
 \|K(t)^{1/2}U(t,s)v\|
 \le \exp\!\left(\frac12\int_s^t\nu(r)\,dr\right)
                 \|K(s)^{1/2}v\|,\qquad v\in\mathcal V.
\end{equation}
For $v\in\mathcal V$ its trajectory is continuous in
$\mathcal H$, bounded in $\mathcal V$, and solves
$i\partial_t U(t,s)v=A(t)U(t,s)v$ in $\mathcal V^*$.
\end{lemma}

\begin{proof}
Fix a reference form $K_*$ and let
$P_n=1_{[1,n]}(K_*)$. These need not have finite rank.
They are contractions in its form and dual norms, converge
strongly there, and their ranges lie in $\mathcal V$.
On $\mathcal H_n=P_n\mathcal H$, the restricted forms
$K_n(t)$ and $A_n(t)=K_n(t)-c(t)I$ are bounded
self-adjoint operators. Uniform equivalence gives their
local uniform operator bounds. The integral form identity
and \eqref{p4f:eq:relativework} give absolute continuity in
operator norm on each such subspace. The integral equation
for the bounded-operator ODE is solved by successive
approximations on short intervals and concatenation; its
Hermitian generator gives a unitary evolution.

Let $u_n(s)=P_nv$. Differentiation of its form energy yields
\[
 \frac d{dt}\langle u_n,K_nu_n\rangle
       =\dot k_t(u_n,u_n).
\]
The two evolution terms cancel because
$A_n=K_n-cI$. Applying \eqref{p4f:eq:relativework}
and Gronwall gives \eqref{p4f:eq:propagation} for $u_n$,
with $P_nv$ at the entrance. In particular $u_n$ is
uniformly bounded in $\mathcal V$. Its derivative,
viewed in the full dual by testing against $P_nw$,
is uniformly bounded in $\mathcal V^*$.

Weak compactness, followed by a diagonal argument on a
countable dense set of form test vectors, gives a limit
$u$ weakly continuous in $\mathcal H$ and weakly bounded
in $\mathcal V$, with $u(s)=v$. In the integrated
equation each test $P_nw$ converges strongly in
$\mathcal V$; uniform boundedness of the forms therefore
gives $i\dot u=A(t)u$ in $\mathcal V^*$.
At every fixed time weak lower semicontinuity of $k_t$
gives the asserted energy bound. This passage uses
conforming compressions of the same forms.

We recall explicitly the norm fact that makes this weak
construction unique. If $w\in L^2(\mathcal V)$ and
$\dot w\in L^2(\mathcal V^*)$, then
\begin{equation}\label{p4f:eq:chain}
 \|w(t)\|^2-\|w(s)\|^2
     =2\operatorname{Re}\int_s^t\langle\dot w(r),w(r)\rangle\,dr,
\end{equation}
and $w$ has a continuous $\mathcal H$ representative.
One obtains this by time smoothing: for smooth
$\mathcal V$-valued functions it is the product rule;
Cauchy--Schwarz in the dual pairing passes the integral
under convergence in $L^2(\mathcal V)$ and
$W^{1,2}(\mathcal V^*)$. The same product rule, integrated
from a time whose norm is bounded by its time average,
bounds the supremum of the $\mathcal H$ norm by these two
space norms. It consequently supplies continuous traces
and justifies the endpoint passage as well.

Apply \eqref{p4f:eq:chain} to the limit and to differences
of solutions. Hermiticity gives
$\operatorname{Re}\langle-iA(t)w,w\rangle=0$, proving norm
conservation and uniqueness before any strong-convergence
claim. Construction backwards from any terminal form
vector gives an inverse evolution. Density extends the
isometries to mutually inverse unitaries on $\mathcal H$.
Uniqueness gives their composition rule and strong
continuity. The displayed energy bound proves invariance
of $\mathcal V$. It also gives weak measurability there;
separability gives the strong measurability needed for
the form-valued integrals below.
\end{proof}

\begin{theorem}[Ordered inverse dressing for dual forcing]
\label{p4f:thm:lift}
Under Assumption~\ref{p4f:ass:forms}, let $q_0\in\mathcal V$ and
$f\in W^{1,1}(0,T;\mathcal V^*)$. Define
\[
 \ell(t)=K(t)^{-1/2}f(t),\qquad
 h(t)=K(t)^{-1/2}\dot f(t),\qquad
 N(t)=\int_0^t\nu(s)\,ds .
\]
Then $i\dot q=A(t)q+f$, $q(0)=q_0$, has a unique solution
in $C([0,T];\mathcal H)\cap L^\infty(0,T;\mathcal V)$.
The inverse lift $g=K^{-1}f$ belongs to
$W^{1,1}(0,T;\mathcal V)$ and obeys the ordered identity
\begin{equation}\label{p4f:eq:inversederivative}
 \dot g=K^{-1}\dot f-K^{-1}\dot K K^{-1}f,\qquad
 K^{1/2}\dot g=h-C_t\ell,\qquad
 C_t=K^{-1/2}\dot K K^{-1/2}.
\end{equation}
For $z=q+g$ and $\mathcal E_z(t)=\|K(t)^{1/2}z(t)\|$,
\begin{align}
 \mathcal E_z(t)
 &\le e^{N(t)/2}\left[
       \mathcal E_z(0)+
       \int_0^t e^{-N(s)/2}
          \{(|c|+\nu)\|\ell\|+\|h\|\}(s)\,ds\right],
       \label{p4f:eq:liftbound}\\
 \|K(t)^{1/2}q(t)\|&\le\mathcal E_z(t)+\|\ell(t)\|.
       \label{p4f:eq:qbound}
\end{align}
The actual entrance term is
$\mathcal E_z(0)=\|K(0)^{1/2}(q_0+K(0)^{-1}f(0))\|$.
\end{theorem}

\begin{proof}
Uniform equivalence and coercivity make each $K(t)$
a bounded isomorphism $\mathcal V\to\mathcal V^*$.
The exact inverse difference identity is
\[
 K(t)^{-1}-K(s)^{-1}
   =-K(t)^{-1}[K(t)-K(s)]K(s)^{-1}.
\]
The relative-work bound makes its norm difference
bounded by a constant times $\int_s^t\nu$.
For a fixed dual vector, the resulting inverse path is
absolutely continuous into the Hilbert space $\mathcal V$.
Strong differentiation of the form identity on fixed
vectors, the inverse difference identity, and uniform
boundedness therefore give
$(K^{-1})'f=-K^{-1}\dot K K^{-1}f$ on fixed vectors.
For completeness, one can choose a common full-measure
set first on a countable dense set of form vectors,
then use the integrable relative bound and Lebesgue
differentiation to extend the difference quotient to
each vector. Approximating the absolutely continuous
path $f$ by simple derivatives proves the product rule
for $K^{-1}f$. This establishes
\eqref{p4f:eq:inversederivative} in $\mathcal V$.
Its second equality is a form-Riesz identification;
no derivative of $K^{-1/2}$ has been taken.

Since $A=K-cI$, direct substitution gives
\begin{equation}\label{p4f:eq:dressed}
 i\dot z=A(t)z+r,\qquad r=cg+i\dot g\in L^1(0,T;\mathcal V).
\end{equation}
Indeed
\[
 \|K^{1/2}r\|\le (|c|+\nu)\|\ell\|+\|h\|.
\]
These terms are integrable: $f$ is bounded in the
fixed dual norm, $\dot f$ is integrable, and the
form norms are uniformly equivalent.

Use Lemma~\ref{p4f:lem:propagation} to define the
form-valued Duhamel integral
\[
 z(t)=U(t,0)(q_0+g(0))
          -i\int_0^tU(t,s)r(s)\,ds.
\]
Its form norm is bounded by the right side of
\eqref{p4f:eq:liftbound}. Fubini in the form-dual
weak equation verifies \eqref{p4f:eq:dressed};
the usual Hilbert-space Duhamel integral is the same
vector, so $z$ is continuous in $\mathcal H$.
Subtracting $g$ gives the asserted solution $q$.
The difference of any two such solutions has derivative
$-iA(t)w$ in $L^\infty(\mathcal V^*)$; the norm chain
rule makes it zero when its entrance is zero.
This proves uniqueness and both inequalities.
\end{proof}

The estimate trades temporal regularity in the form dual
for spatial first-form control. It does not require
$f\in\mathcal V$ or $Af\in\mathcal H$.
For an actual coupling $f=L(t)p(t)$, sufficient explicit
input contracts are
\begin{align}
 \|K^{-1/2}LF^{-1}\|&\le b(t),&
 \|K^{-1/2}\dot L F^{-1}\|&\le b_1(t),\notag\\
 \|\ell\|&\le b\|Fp\|,&
 \|h\|&\le b_1\|Fp\|+b\|F\dot p\|.
 \label{p4f:eq:inputcontract}
\end{align}
Here $F\ge I$ is a fixed input graph and the products
and derivatives must hold on their stated domains.
The last quantity is a property of the coupled input,
not of a frozen occupation. A changing $F$ adds its
own derivative or commutator.

\begin{corollary}[Retained internal right clock]
\label{p4f:cor:clock}
For a Hilbert--Schmidt factor satisfying
$i\dot q=Aq-qD+f$, with fixed bounded Hermitian
right clock $D$, the same argument applies with
\[
 r=K^{-1}\{i\dot f+cf+fD-i\dot K K^{-1}f\}.
\]
In \eqref{p4f:eq:liftbound} its inhomogeneous norm can
be bounded by
$\|K^{-1/2}(i\dot f+cf+fD)\|_{\rm HS}
 +\nu\|\ell\|_{\rm HS}$.
\end{corollary}
\begin{proof}
With $z=q+g$, substitution gives
$i\dot z=Az-zD+cg+gD+i\dot g$.
Right multiplication by the clock unitary removes
$-zD$ and preserves Hilbert--Schmidt norms, including
the physical left form norm. Apply the preceding
Duhamel estimate and substitute \eqref{p4f:eq:inversederivative}.
\end{proof}

Only a bounded physical-coordinate-independent internal
clock is covered by this corollary. An unbounded right
source clock needs its own graph contract. A positional
purifier cannot be changed into an internal index just
because the Hilbert-space expressions have the same size.
The scalar $c$ is not free: it appears in the full
combination $i\dot f+cf+fD$, while $\dot c$ already belongs
to $\dot K$. Phases may cancel a term only through that
complete physical combination.

\subsection{Two sharp boundaries of the temporal argument}

\begin{example}[Continuous dual forcing need not give a wave]
\label{p4f:ex:roughforcing}
On $\ell^2(n\ge2)$ let $A_n=n^2$, $K_n=n^2+1$, and
$\ell_n(t)=n^{-1}e^{-in^2t}$. Dominated convergence
makes $\ell(t)$ continuous in $\ell^2$; hence
$f=K^{1/2}\ell$ is continuous in $\mathcal V^*$.
Its zero-initial causal solution has components
\[
 q_n(t)=-it\,\frac{\sqrt{n^2+1}}n e^{-in^2t}.
\]
For every $t>0$ these are not square summable.
Thus a bounded form-dual coupling and continuous
input alone do not give a physical Hilbert-space
response. The time-derivative hypothesis in
Theorem~\ref{p4f:thm:lift} fails.
\end{example}
\begin{proof}
The dual group is diagonal, so substitution in the
causal integral gives the displayed formula.
Its component magnitude tends to $t$. On the other
hand the derivative of $\ell_n$ has magnitude $n$,
so it cannot supply the required integrable dual
derivative of $f$.
\end{proof}

\begin{proposition}[Why an inverse square root is not the lift]
\label{p4f:prop:sqrtwarning}
The relative form-speed condition does not give a
dimension-independent bound on
$K^{1/2}(K^{-1/2})'$. It does give an ordered bounded
inverse derivative as in \eqref{p4f:eq:inversederivative}.
\end{proposition}
\begin{proof}
In a finite eigenbasis, write
$\dot K=K^{1/2}CK^{1/2}$ and $K_{ii}=\lambda_i>0$.
Differentiate $(K^{1/2})^2=K$ and then its inverse;
the Sylvester equation gives
\begin{align*}
 (K^{-1/2})'_{ij}
   &=-\frac{C_{ij}}{\sqrt{\lambda_i}+\sqrt{\lambda_j}},\\
 [K^{1/2}(K^{-1/2})']_{ij}
   &=-\frac{\sqrt{\lambda_i}C_{ij}}
                     {\sqrt{\lambda_i}+\sqrt{\lambda_j}}.
\end{align*}
The first kernel equals
$-\int_0^\infty e^{-sK^{1/2}}C e^{-sK^{1/2}}\,ds$,
so its norm is at most
$\|C\|/(2\sqrt{\lambda_{\min}})$. The warning concerns
the second, weighted derivative.

Take $\lambda_j=4^j$ and the Hermitian matrix
$C_{ij}=i/[\pi(i-j)]$ for $i\ne j$, $C_{ii}=0$,
on indices $1,\ldots,N$. Its norm is at most one:
it is a compression of the convolution operator on
$\ell^2(\mathbb Z)$ whose Fourier multiplier is,
up to its sign, $(\theta-\pi)/\pi$ on $(0,2\pi)$.
This follows by integrating the Fourier series of
the sawtooth, or by Abel summing
$\sum_{d\ne0}e^{id\theta}/d=i(\pi-\theta)$.
For the unit vector $v=N^{-1/2}(1,\ldots,1)$, pairing
the two off-diagonal entries at separation $d$ gives
\[
 \left|\left\langle v,
       K^{1/2}(K^{-1/2})'v\right\rangle\right|
 =\frac1{\pi N}\sum_{d=1}^{N-1}
       \frac{(N-d)\tanh(d\log2/2)}d .
\]
For $2\le d\le N/2$, the hyperbolic tangent is at
least $3/5$ and $(N-d)/N\ge1/2$. The harmonic sum
diverges, proving the assertion. These are actual
positive paths
$K(s)=K(0)^{1/2}(I+sC)K(0)^{1/2}$, $|s|\leq1/2$.
Since $K(0)\geq4I$, these obey $K(s)\geq2I$; their relative
speed is bounded by $(1-|s|)^{-1}\leq2$.
The ordered full inverse in the theorem avoids this
weighted-square-root obstruction.
\end{proof}

\section{A finite spectral window with a varying coherent input}
\label{p4f:sec:window}

We next fix the complete self-adjoint complement $A$.
Its Hilbert space may contain continua and all quantum
source coordinates. Let $\mathcal U$ be the input Hilbert
space, $B:\mathcal U\to\mathcal H$ bounded, and
$a\in L^2(0,T;\mathcal U)$. The zero-initial response is
\begin{equation}\label{p4f:eq:boundedresponse}
 q_B(t)=-i\int_0^t e^{-iA(t-s)}B a(s)\,ds .
\end{equation}
A spectral sector can be included by replacing $B$ with
$CB$ for a spectral projection $C$ of $A$.
The input may change direction arbitrarily in $\mathcal U$.

\begin{lemma}[An operator spectral measure on one interval]
\label{p4f:lem:interval}
Let $I=[\alpha,\alpha+l)$, $l>0$, and suppose
$B^*E_A(I)B\le mI_{\mathcal U}$. For
$v\in H^1(I;\mathcal U)$, the spectral integral
$\mathcal T_Iv=\int_I E_A(d\lambda)Bv(\lambda)$ satisfies
\begin{equation}\label{p4f:eq:interval}
 \|\mathcal T_Iv\|^2
 \le2m\left\{l^{-1}\|v\|_{L^2(I)}^2+
                         l\|v'\|_{L^2(I)}^2\right\}.
\end{equation}
\end{lemma}

\begin{proof}
Let $\bar v=l^{-1}\int_Iv$ be the Bochner mean.
The fundamental theorem for Hilbert-valued $H^1$ functions
gives
\[
 v(\lambda)=\bar v+\int_I k(\lambda,\xi)v'(\xi)\,d\xi,
 \qquad
 k(\lambda,\xi)=1_{\xi<\lambda}
                    -\frac{\alpha+l-\xi}{l},\qquad |k|\le1.
\]
For every scalar $b$ supported in $I$ with $|b|\le1$,
\[
 \|b(A)E_A(I)B\|^2
 =\|B^*E_A(I)|b(A)|^2B\|\le m.
\]
Consequently the spectral integral has the precise order
\[
 \mathcal T_Iv=E_A(I)B\bar v+
     \int_I k(A,\xi)E_A(I)Bv'(\xi)\,d\xi.
\]
For smooth $v$ this follows by its displayed integral
representation and bounded spectral calculus. The same
formula defines a continuous extension to $H^1$.
It gives
\[
 \|\mathcal T_Iv\|
 \le\sqrt m\left\{l^{-1/2}\|v\|_{L^2(I)}
                         +l^{1/2}\|v'\|_{L^2(I)}\right\}.
\]
Squaring and using $(x+y)^2\le2x^2+2y^2$ proves
\eqref{p4f:eq:interval}. No common eigenvector of the
input-valued function or simultaneous diagonalization
of the operator measure was assumed.
\end{proof}

\begin{theorem}[Finite-imaginary-part response window]
\label{p4f:thm:window}
For $\eta>0$ define the positive operator
\begin{equation}\label{p4f:eq:poisson}
 M_\eta(E)=\frac1\pi B^*
             \frac{\eta}{(A-E)^2+\eta^2}B .
\end{equation}
Suppose $M_\eta(E)\le s_\eta I$ for all real $E$,
or at the centers of a partition into intervals of
length $2\eta$ covering the relevant spectrum.
Then, for $0\le t\le T$,
\begin{equation}\label{p4f:eq:window}
 \|q_B(t)\|^2\le
 4\pi^2s_\eta(1+\eta^2t^2)
                   \int_0^t\|a(s)\|^2\,ds .
\end{equation}
Atoms, singular spectral measures and thresholds are allowed.
The coefficient is an operator ceiling for the complete
input space, not a scalar response for one incident column.
\end{theorem}

\begin{proof}
On $I_j=[E_j-\eta,E_j+\eta)$ the Poisson kernel in
\eqref{p4f:eq:poisson} is at least $1/(2\pi\eta)$.
Hence
$B^*E_A(I_j)B\le2\pi\eta s_\eta I$.
Set
\[
 v_t(\lambda)=\int_0^t
                  e^{i\lambda(s-t/2)}a(s)\,ds .
\]
This Bochner Fourier transform is in
$H^1(\mathbb R;\mathcal U)$, even though no derivative
of $a$ was required. Hilbert-valued Plancherel gives
\begin{align*}
 \|v_t\|_{L^2(\mathbb R)}^2
     &=2\pi\int_0^t\|a(s)\|^2\,ds,\\
 \|v_t'\|_{L^2(\mathbb R)}^2
     &=2\pi\int_0^t(s-t/2)^2\|a(s)\|^2\,ds
       \le\frac{t^2}{4}\|v_t\|_{L^2(\mathbb R)}^2.
\end{align*}
The full spectral integral of $v_t$ differs from
$q_B(t)$ only by $-i e^{-iAt/2}$. This identity follows
directly by Fubini from \eqref{p4f:eq:boundedresponse}.
The outputs $\mathcal T_{I_j}v_t$ lie in mutually
orthogonal spectral subspaces. Apply
Lemma~\ref{p4f:lem:interval} with
$l=2\eta$, $m=2\pi\eta s_\eta$ and sum:
\[
 \|q_B(t)\|^2
 \le 2\pi s_\eta\|v_t\|_2^2+
                    8\pi\eta^2s_\eta\|v_t'\|_2^2 .
\]
The Plancherel identities give
\eqref{p4f:eq:window}. Orthogonality and monotone
summation justify infinitely many intervals without
a dimension or number-of-channels factor.
\end{proof}

The interval derivative term is substantive. A ceiling
on $B^*E_A(I)B$ by itself cannot bound a spectral integral
of a changing vector by a dimension-free multiple of
$\sup_\lambda\|v(\lambda)\|^2$. The time-centered Fourier
representation is what supplies the derivative in
this theorem. It is a derivative in spectral energy,
not a time-regularity premise on the actual input.

For intervals of half-width $\delta$ instead of $\eta$,
the same proof gives
\begin{equation}\label{p4f:eq:otherwindow}
 \|q_B(t)\|^2\le
 2\pi^2s_\eta\frac{\eta^2+\delta^2}{\eta}
           (\delta^{-1}+\delta t^2)
                      \int_0^t\|a(s)\|^2\,ds.
\end{equation}
Indeed the interval measure costs
$m=\pi(\eta^2+\delta^2)s_\eta/\eta$ and its length is
$2\delta$. Thus the observation window can be chosen
to match an earned response table.

If $B=LF^{-1}$, use $a=Fp$ with the actual input.
For $F=I$ and a normalized full block evolution,
$\int_0^t\|p(s)\|^2ds\le t$; choosing $\eta=1/T$
gives the uniform ceiling $8\pi^2s_{1/T}T$.
If $F$ is dimensionless and $L,A$ have inverse-time
units, $s_\eta$ has inverse-time units, so this norm
bound is dimensionless.
An atom of response weight $w$ at $\lambda_0$ has
$M_\eta(\lambda_0)=w/(\pi\eta)$, whereas a resonant
constant input can give response $wt^2$.
The theorem retains that pole through $s_\eta$; it
does not convert it into a decay rate or a gap.

\subsection{Unbounded form coupling: near response and far dressing}

\begin{theorem}[Coherent near/far response in the physical form]
\label{p4f:thm:hybrid}
Let $A$ be fixed self-adjoint, $K=A+c\ge I$, and
$F\ge I$ a fixed positive input graph. Assume
\[
 B=K^{-1/2}LF^{-1}:\mathcal U\longrightarrow\mathcal H
 \quad\hbox{is bounded}.
\]
Thus $LF^{-1}=K^{1/2}B$ is a bounded map into
$\mathcal V^*$, possibly not into $\mathcal H$.
Assume $Fp\in L^2(0,T;\mathcal U)$ for the near-response statement.
For $J_0=[E_0-\Lambda,E_0+\Lambda]$, $\Lambda>0$,
let $C=1_{J_0}(A)$. For $j=0,1$ set
\[
 B_j=K^{(j+1)/2}CB,\qquad
 M_{j,\eta}(E)=\frac1\pi B^*C K^{j+1}
                    \frac{\eta}{(A-E)^2+\eta^2}CB .
\]
These $B_j$ are bounded and $M_{j,\eta}$ are positive
operator spectral responses. If
$M_{j,\eta}(E)\le s_{j,\eta}I$ at the required centers,
the zero-initial near response satisfies
\begin{equation}\label{p4f:eq:near}
 \|K^{j/2}q_{\rm near}(t)\|^2
 \le4\pi^2s_{j,\eta}(1+\eta^2t^2)
                          \int_0^t\|Fp(s)\|^2\,ds .
\end{equation}
Let
\[
 R_j=K^{(j+1)/2}(A-E_0)^{-1}(1-C)B .
\]
If $a_0(s)=e^{iE_0s}Fp(s)\in W^{1,1}(0,t;\mathcal U)$,
the far response satisfies
\begin{align}
 \|K^{j/2}q_{\rm far}(t)\|
 &\le\|R_j\|
   \left\{\|a_0(t)\|+\|a_0(0)\|
                       +\int_0^t\|\dot a_0(s)\|\,ds\right\},
                       \label{p4f:eq:far}\\
 \|R_0\|&\le
       \frac{\sqrt{\Lambda+|E_0+c|}}{\Lambda}\|B\|,
 &\|R_1\|&\le
       \left(1+\frac{|E_0+c|}{\Lambda}\right)\|B\|.
                       \label{p4f:eq:farcoeff}
\end{align}
An original $q_0\in\mathcal V$ contributes the separate
homogeneous term $e^{-iAt}q_0$, with its original
form norm. It is not part of \eqref{p4f:eq:near}.
\end{theorem}

\begin{proof}
The causal integral is first defined in $\mathcal V^*$,
where $e^{-iAt}$ is a unitary group in the $K$ dual
norm. On the bounded spectral interval, multiplication
by $K^{j/2}C$ turns its forcing into $B_jFp$.
Apply Theorem~\ref{p4f:thm:window} to this bounded
coupling to obtain \eqref{p4f:eq:near}.

On the far sector put $D=A-E_0$. Commuting only
functions of this same self-adjoint $A$, integration
by parts gives the exact identity
\begin{align}
 K^{j/2}q_{\rm far}(t)
 =e^{-iE_0t}\bigg\{
 &e^{-iDt}R_j a_0(0)-R_j a_0(t)\notag\\
 &+\int_0^t e^{-iD(t-s)}R_j\dot a_0(s)\,ds\bigg\}.
 \label{p4f:eq:faridentity}
\end{align}
For an unbounded coupling, prove it first on bounded
spectral cutoffs. The bounds below pass it to the
form-dual causal solution and also show that its
right side is a physical $K^{j/2}$ vector.
Taking norms proves \eqref{p4f:eq:far}.
The two displayed endpoints are the virtual
dressings; neither can be dropped.

For $x=|\lambda-E_0|\ge\Lambda$ and
$\lambda\in\sigma(A)$, positivity and the triangle
inequality give
$0<\lambda+c\le x+|E_0+c|$. Therefore
\[
 \frac{\sqrt{\lambda+c}}{|\lambda-E_0|}
 \le\frac{\sqrt{\Lambda+|E_0+c|}}{\Lambda},
 \qquad
 \frac{\lambda+c}{|\lambda-E_0|}
 \le1+\frac{|E_0+c|}{\Lambda}.
 \]
The spectral theorem proves \eqref{p4f:eq:farcoeff}
and the limiting integration-by-parts assertion.
The homogeneous term follows by linearity; it has
unchanged $K$ norm because $K=A+c$ commutes with $A$.
\end{proof}

One may use $M_{1,\eta}\le k_{J_0}M_{0,\eta}$, where
$k_{J_0}=\sup_{\lambda\in J_0\cap\sigma(A)}(\lambda+c)$,
if that additional factor is charged. A smaller
first-form response coefficient is a separate spectral
calculation. For several coherent couplings
$L(t)=\sum_\alpha l_\alpha(t)L_\alpha$, apply the theorem
to the stacked bounded map with input
$(l_1Fp,\ldots,l_mFp)$. Its full operator measure retains
all off-diagonal responses. Summing independent scalar
channel rates would be a different statement.

The near estimate requires an $L^2$ input but no
input time derivative. The far estimate requires the
displayed derivative and endpoint dressings. Neither
frozen spectral statement applies automatically to
a changing $A(t)$. One must instead use
Theorem~\ref{p4f:thm:lift}, or prove a comparison with
a fixed parent and charge its actual forcing and
form-dual derivative.

\section{Ordered resolvents for an unbounded closed-form interaction}
\label{p4f:sec:ordered}

\begin{theorem}[Closed-form factor transfer and physical weights]
\label{p4f:thm:ordered}
Let $A_0$ and $A=A_0+W^*JW$ be self-adjoint operators
defined by closed semibounded forms on the same
$\mathcal V$, with equivalent positive shifted form norms.
Here $W:\mathcal V\to\mathcal H_{\rm aux}$ is bounded
and $J=J^*$ is bounded on $\mathcal H_{\rm aux}$.
The adjoint $W^*$ maps $\mathcal H_{\rm aux}$ to
$\mathcal V^*$. Actual form closure is an assumption;
a formal factorization alone does not supply it.

For $\zeta=E+i\eta$, $\eta>0$, set
\[
 R_0=(A_0-\zeta)^{-1},\quad R=(A-\zeta)^{-1},\quad
 G_0=WR_0W^*,\quad G=WRW^* .
\]
All products are bounded between the indicated form
spaces, and
\begin{equation}\label{p4f:eq:ordered}
 T=(I+JG_0)^{-1}=I-JG,\qquad
 RW^*=R_0W^*T,\qquad G=G_0T .
\end{equation}
In particular
\begin{align}
 \operatorname{Im}G
   &=T^*(\operatorname{Im}G_0)T
     =\eta(RW^*)^*(RW^*), \label{p4f:eq:imag}\\
 (RW^*)^*T_{\rm ph}(RW^*)
   &=T^*\big[(R_0W^*)^*T_{\rm ph}(R_0W^*)\big]T .
       \label{p4f:eq:physicalsandwich}
\end{align}
The second formula is an equality of bounded
auxiliary-space forms for any positive closed physical
form $T_{\rm ph}$ with
$\mathcal V\subset\operatorname{Dom}T_{\rm ph}^{1/2}$
continuously. No commutation with $A_0,A,W$ or $J$
is required.
\end{theorem}

\begin{proof}
Choose $K_0=A_0+c_0\ge I$. The spectral theorem gives
the exact identity
\begin{equation}\label{p4f:eq:spectralsup}
 \|K_0^{1/2}R_0K_0^{1/2}\|
  =\sup_{\lambda\in\sigma(A_0)}
                  \frac{\lambda+c_0}{|\lambda-\zeta|}<\infty .
\end{equation}
Thus $R_0$ extends boundedly $\mathcal V^*\to\mathcal V$.
The same argument with a positive shift of $A$ and
equivalence of form norms gives this extension for $R$.
Their form inverse identities therefore give both
ordered resolvent identities
\[
 R-R_0=-R_0W^*JWR=-RW^*JWR_0.
\]
Sandwiching by $W,W^*$ yields
$G-G_0=-G_0JG=-GJG_0$.
It follows by multiplication in the displayed order
that both products of $I+JG_0$ and $I-JG$ equal $I$.
This proves invertibility and the formula for $T$.
The first resolvent identity now gives
$RW^*=R_0W^*(I-JG)$, proving \eqref{p4f:eq:ordered}.

For a self-adjoint resolvent,
$\operatorname{Im}R=\eta R^*R$. The identity extends
to dual inputs by continuity: both sides define
bounded sesquilinear forms on $\mathcal V^*$, and
Hilbert inputs are dense there. Sandwiching with
$W^*$ proves the second equality in
\eqref{p4f:eq:imag}. Substitution of
$RW^*=R_0W^*T$ proves its first equality.
The same substitution into the quadratic form of
$T_{\rm ph}$ gives \eqref{p4f:eq:physicalsandwich}.
Continuity of that form on $\mathcal V$ justifies
each product even when $T_{\rm ph}$ is unbounded
on $\mathcal H$.
\end{proof}

The resolvent inverse in this theorem exists off the
real axis because the closed self-adjoint parents
exist. A Neumann criterion is sufficient for a
quantitative bound, not necessary for the identity:
\begin{equation}\label{p4f:eq:neumann}
 \|JG_0\|\le\theta<1\quad\Longrightarrow\quad
                       \|T\|\le(1-\theta)^{-1}.
\end{equation}
If $LF^{-1}=W^*b$ with bounded $b$, an earned
$\operatorname{Im}G_0/\pi\le s_{0,\eta}I$ implies
\[
 \frac1\pi\operatorname{Im}
       [(LF^{-1})^*R(LF^{-1})]
 \le \frac{\|b\|^2s_{0,\eta}}{(1-\theta)^2}I
\]
when its full form-dual sandwich is used. Likewise
\eqref{p4f:eq:physicalsandwich} transfers an earned
physical-energy resolvent bound. Absorptive response
alone does not replace that physical-energy bound.
For a form-dual coupling this is first a resolvent
sandwich: one uses the bounded near factors of
Theorem~\ref{p4f:thm:hybrid}, or the ordered temporal
estimate below, rather than treating $LF^{-1}$ as
a bounded Hilbert-space coupling without proof.

Equality in \eqref{p4f:eq:spectralsup} is over the
actual spectrum. Replacing it by a whole half-line
above a spectral floor gives an upper bound, not
in general an equality. Spectral gaps matter for
that distinction, whereas the shift that makes
$K_0$ positive does not create such gaps.
At a real eigenvalue no inverse has been asserted.
As $\eta$ decreases, genuine poles and thresholds
may make both the baseline response and the
quantitative inverse ceiling large.

\section{Damped temporal response with a noncommuting physical form}
\label{p4f:sec:damped}

The physical derivative form generally does not commute
with the interacting complement. It may therefore be
placed outside an actual resolvent, as below, but cannot
be inserted into a scalar spectral measure as though
it commuted with every spectral projection.
This use of a resolvent and a time integral is related
to the Kato smoothing framework
\cite{Kato1966,DAncona2014}. The statement here is at
one earned $\eta>0$; it does not assert the uniform
imaginary-part bounds required by a global smoothing
theorem.

\begin{theorem}[Ordered damped response]
\label{p4f:thm:damped}
Let $A$ be fixed self-adjoint, $K=A+c\ge I$, and
$\mathcal L=LF^{-1}:\mathcal U\to\mathcal V^*$ bounded.
Let $T_{\rm ph}\ge I$ be a positive closed form whose
domain contains $\mathcal V$ continuously.
For fixed $\eta>0$ suppose the actual ordered sandwich
satisfies the measurable bound
\begin{equation}\label{p4f:eq:orderedbound}
 \|T_{\rm ph}^{1/2}(A-E-i\eta)^{-1}\mathcal L\|^2
                         \le S_\eta(E).
\end{equation}
For $a\in L^2(0,T;\mathcal U)$, extend the forcing by
zero after $T$, and let $q$ be its causal zero-initial
form-dual response. Define
$a_\eta(s)=e^{-\eta s}a(s)1_{[0,T]}(s)$ and use
the unitary Fourier convention
$\widehat u(E)=(2\pi)^{-1/2}\int_{\mathbb R}e^{iEt}u(t)\,dt$.
If the right side is finite, then
\begin{equation}\label{p4f:eq:damped}
 \int_0^\infty e^{-2\eta t}
                      \|T_{\rm ph}^{1/2}q(t)\|^2\,dt
 \le \int_{\mathbb R}S_\eta(E)
                              \|\widehat a_\eta(E)\|^2\,dE .
\end{equation}
In particular a uniform $S_\eta(E)\le S_\eta$ gives
the upper bound $S_\eta\|a_\eta\|_2^2$. A bound
$S_\eta(E)\le S_0+S_1E^2$ instead gives
\begin{equation}\label{p4f:eq:temporalprice}
 S_0\|a_\eta\|_2^2+S_1\|\dot a_\eta\|_2^2,
\end{equation}
provided the actual zero extension $a_\eta$ belongs
to $H^1(\mathbb R;\mathcal U)$.
\end{theorem}

\begin{proof}
The fixed group $e^{-iAt}$ is unitary also on
the $K$ form dual. The causal integral
\[
 q(t)=-i\int_0^{\min(t,T)}
                    e^{-iA(t-s)}\mathcal L a(s)\,ds
\]
is therefore a continuous $\mathcal V^*$ vector.
It is zero for negative times and bounded in that
dual norm for all positive times. Its damped
extension $q_\eta=e^{-\eta t}q\,1_{[0,\infty)}$
belongs to $L^2(\mathbb R;\mathcal V^*)$.
Its value at zero is zero, so differentiation
creates no entrance delta distribution.
Fourier transformation of the causal equation,
or Fubini applied to its damped convolution, gives
the exact form-dual identity
\begin{equation}\label{p4f:eq:fourierresponse}
 \widehat q_\eta(E)
       =-(A-E-i\eta)^{-1}\mathcal L\widehat a_\eta(E).
\end{equation}

For almost every $E$ the right side lies in
$\mathcal V$, hence in
$\operatorname{Dom}T_{\rm ph}^{1/2}$. The assumed
finite integral and \eqref{p4f:eq:orderedbound}
place it in the $L^2$ graph of
$T_{\rm ph}^{1/2}$. Coercivity $T_{\rm ph}\ge I$
also places it in physical $L^2(\mathcal H)$.
Fourier transformation commutes with this closed
spatial operator: to see this, replace it by its
bounded spectral truncations, use Hilbert-valued
Plancherel there, and pass by its closed graph
and monotone convergence of the graph norms.
Thus inverse Fourier transformation identifies
this same dual solution as a physical
$L^2(\operatorname{Dom}T_{\rm ph}^{1/2})$ response.
Plancherel and \eqref{p4f:eq:orderedbound} prove
\eqref{p4f:eq:damped}. In this argument the
truncations are of $T_{\rm ph}$ for the Fourier
graph identity; none is commuted through $A$
or its resolvent.

For a uniform coefficient use Plancherel on
$a_\eta$. For the quadratic coefficient use
$\|E\widehat a_\eta\|_2=\|\dot a_\eta\|_2$.
That equality requires the stated whole-line
$H^1$ condition and proves \eqref{p4f:eq:temporalprice}.
\end{proof}

The theorem gives an integrated first-form response,
not an all-times pointwise form bound. If the input
is $H^1$ inside $[0,T]$ but has nonzero endpoint
values, its zero extension has delta derivatives
and is not $H^1(\mathbb R)$. One must keep the
endpoint dressings, as in
\eqref{p4f:eq:faridentity} or
Theorem~\ref{p4f:thm:lift}, or prove a physical
pulse collar that removes those jumps.
The endpoint problem cannot be repaired by simply
omitting the positive frequency-weighted term.

The physical finite-interval graph norm obeys
\begin{equation}\label{p4f:eq:undamping}
 \|T_{\rm ph}^{1/2}q\|_{L^2(0,T)}
 \le e^{\eta T}
       \left(\int_{\mathbb R}S_\eta(E)
                          \|\widehat a_\eta(E)\|^2\,dE\right)^{1/2}.
\end{equation}
If $q_0\ne0$, its homogeneous response must be
added separately. For example, an earned
$T_{\rm ph}\le\kappa K$ and $q_0\in\mathcal V$
give
\[
 \left(\int_0^\infty e^{-2\eta t}
       \|T_{\rm ph}^{1/2}e^{-iAt}q_0\|^2dt\right)^{1/2}
 \le\sqrt{\frac{\kappa}{2\eta}}\|K^{1/2}q_0\|.
 \]
Add this to the square root of the forced bound;
the square of the sum, not the sum of squares,
is a generally valid full-response bound.
For a noncoercive derivative form, replace it by
$I+T_{\rm ph}$ or establish a separate Hilbert-norm
sandwich. Temporal damping has neither projected
out a bound pole nor supplied a missing source graph.

\section{Whole-form spatial and source derivatives}
\label{p4f:sec:derivatives}

Temporal smoothing and spatial differentiation
make different demands. The preceding theorems
do not require a spatial derivative of the forcing
when their temporal or resolvent assumptions are
earned. A physical source derivative, however,
must act on the complete coupled form and the
actual input; replacing that coordinate by a
classical label changes the problem.

\begin{proposition}[Ordered whole-form derivative]
\label{p4f:prop:derivatives}
Let $A$ and $\mathcal L=LF^{-1}$ be as in
Theorem~\ref{p4f:thm:damped}. Let $\delta_i$ be
either a parameter derivative on the common
form domain, or a closed physical derivation
$\delta_iM=D_i^QM-MD_i^P$.
For a parameter family, require strong difference-quotient convergence
of $A$ in $\mathcal L(\mathcal V,\mathcal V^*)$ and of $\mathcal L$
in $\mathcal L(\mathcal U,\mathcal V^*)$, with locally uniform bounds
on these quotients and uniformly equivalent nearby form norms.
For a physical derivation, require domain-preserving translations or
local gauges for which the conjugated families satisfy these same
conditions and realize the displayed derivatives on a common core.
Suppose the actual whole-form sandwiches are bounded:
\begin{equation}\label{p4f:eq:derivativecontracts}
 \|K^{-1/2}(\delta_iA)K^{-1/2}\|\le\alpha_i,\qquad
 \|K^{-1/2}\delta_i\mathcal L\|\le\beta_i .
\end{equation}
For $\zeta=E+i\eta$, $\eta>0$, write
$R_\zeta=(A-\zeta)^{-1}$ and
$U_\zeta=K^{1/2}R_\zeta\mathcal L$. Then
\begin{align}
 \delta_i(R_\zeta\mathcal L)
   &=R_\zeta\{\delta_i\mathcal L
                    -(\delta_iA)R_\zeta\mathcal L\},
                    \label{p4f:eq:derivativeidentity}\\
 \|K^{1/2}\delta_i(R_\zeta\mathcal L)\|
   &\le J_\zeta
          \{\beta_i+\alpha_i\|U_\zeta\|\},\qquad
 J_\zeta=\sup_{\lambda\in\sigma(A)}
                      \frac{\lambda+c}{|\lambda-\zeta|}.
                    \label{p4f:eq:derivativebound}
\end{align}
If $\zeta$ itself varies, add
$(\delta_i\zeta)R_\zeta^2\mathcal L$ to
\eqref{p4f:eq:derivativeidentity}.
\end{proposition}

\begin{proof}
For a parameter family, the form resolvent
difference identity is
\[
 R_{\zeta,\epsilon}-R_\zeta
   =-R_{\zeta,\epsilon}(A_\epsilon-A)R_\zeta.
\]
It is a bounded identity from $\mathcal V^*$
to $\mathcal V$. Strong differentiation on
fixed vectors, uniform local form bounds and
the given difference-quotient hypothesis yield
$\delta_iR_\zeta=-R_\zeta(\delta_iA)R_\zeta$.
Leibniz gives \eqref{p4f:eq:derivativeidentity}.
For a physical derivation, conjugate by the
domain-preserving translations/gauges first
and apply the same argument; differentiation
on the common core is the commutator formula.
Strong form convergence and closedness of the
physical derivative pass it to its admitted domain.

Insert a form Riesz map without changing order:
\[
 K^{1/2}R_\zeta(\delta_iA)R_\zeta\mathcal L
 =\big[K^{1/2}R_\zeta K^{1/2}\big]
  \big[K^{-1/2}(\delta_iA)K^{-1/2}\big]U_\zeta .
\]
The first factor has norm $J_\zeta$ by the
spectral theorem. The analogous factorization
of the $\delta_i\mathcal L$ term proves
\eqref{p4f:eq:derivativebound}. Differentiating
$A-\zeta$ when $\zeta$ varies gives the stated
extra term with its positive sign.
\end{proof}

For a literal input vector $a$, the corresponding
physical derivative remains
\begin{equation}\label{p4f:eq:inputspatial}
 D_i^Q(R_\zeta\mathcal L a)
   =[\delta_i(R_\zeta\mathcal L)]a
                        +R_\zeta\mathcal L D_i^Pa .
\end{equation}
Its second term cannot be erased by a response
norm estimate. The graph and derivative domain
of $a$ must be included. When $F$ is covariantly
constant, $\delta_i\mathcal L=(\delta_iL)F^{-1}$;
otherwise
\[
 \delta_i\mathcal L
    =(\delta_iL)F^{-1}
                -LF^{-1}(\delta_iF)F^{-1}
\]
on the admitted graph, with an additional bound
for that second term. Gauge curvature, moving
frames, source kinetic terms and mobility
commutators all belong to the whole forms in
\eqref{p4f:eq:derivativecontracts}.
Differentiating a singular factor $W$ is unnecessary
and may be invalid even when the derivative of
the whole form $W^*JW$ is bounded.

For example, in three relative dimensions the
Hardy inequality gives the legitimate whole
Coulomb force form
\[
 |\langle u,\partial_{R_i}(g/|r-R|)v\rangle|
 \le |g|\|u/|r-R|\|\|v/|r-R|\|
 \le4|g|\|\nabla_r u\|\|\nabla_r v\|.
\]
With an earned $K\ge\lambda_{\rm rel}(-\Delta_r)$,
its relative coefficient is at most
$4|g|/\lambda_{\rm rel}$. The coefficient is
the true relative-mass kinetic coefficient;
source-dependent center maps add their actual
derivatives. This is a form estimate, not an
$L^2$ operator bound on the force. Strong vector
continuity or the stated difference-quotient
contract suffices; operator-norm continuity
under translations is not assumed.

All statements above keep a fixed physical block
chart. For a varying projector one must include
its domain-preserving transport and the moving
frame term $-iU^*\dot U$. A nonlocal Hilbert-space
band unitary is not automatically an isometry
of pointwise physical currents. Likewise a
moving hard-core boundary requires a proved
domain transport before it enters a
fixed-domain force identity.

\section{Spatial force and temporal regularity as distinct current resources}
\label{p4s:sec:forms}

The form-dual estimates above and the following estimates solve different
input problems. A residual with a spatial form norm can be propagated
using actual forces without paying an irrelevant constant carrier
energy. A residual with only a Hilbert norm can instead use temporal
regularity and a common strong domain. Neither hypothesis implies the
other for singular interactions.

\subsection{A force estimate with covariant drift}

\begin{theorem}[Forced positive form with complete first-order work]
\label{p4s:thm:force}
Use \eqref{p4s:eq:parent} and
$K=-\sum_i\lambda_iD_i^2+U\geq I$, with scalar $U\geq1$.
Assume a common form domain, the true unitary propagator, and a
form-stable approximation justifying the commutators below. Suppose
$K'\leq a_KK$ as forms and, for form vectors $v$,
\begin{align*}
 \left(\sum_i\|\sqrt{\lambda_i}(\nabla_iV)v\|^2\right)^{1/2}
       &\leq g_0\|v\|+g_1\|K^{1/2}v\|,\\
 \|B\|&\leq L,\qquad
 B_{ji}=\sqrt{\lambda_j/\lambda_i}\,\nabla_j a_i,\\
 \left(\sum_j\|\sqrt{\lambda_j}A_jv\|^2\right)^{1/2}
       &\leq G_0\|v\|+G_1\|K^{1/2}v\|,\\
 \left\|U^{-1/2}\sum_i a_i(\partial_iU)v\right\|
       &\leq u_0\|v\|+u_1\|K^{1/2}v\|,
\end{align*}
where $B$ acts on the coordinate direct sum, and
\[
 [D_j,D_i]=-i\Omega_{ji},\qquad
 A_j=\sum_i\left(-ia_i\Omega_{ji}
                          +\tfrac12\nabla_j\nabla_i a_i\right).
\]
All scalar coefficients are nonnegative and integrable.
If $iE'=HE+r$, $E(0)$ is in the form domain and
$r\in L^1([0,T];\operatorname{Dom}K^{1/2})$, then
\begin{align}
 e(t)&\leq e_0+\int_0^t\|r(s)\|\,ds=:E_*(t),\nonumber\\
 Q(t)&\leq e^{B_*(t)}
  \left[Q_0+\int_0^t e^{-B_*(s)}
       \left(\beta(s)E_*(s)+\|K(s)^{1/2}r(s)\|\right)ds\right],
       \label{p4s:eq:force}\\
 B_*(t)&=\int_0^t\alpha(s)\,ds,\quad
 \alpha=\tfrac12a_K+g_1+L+G_1+\tfrac12u_1,\quad
 \beta=g_0+G_0+\tfrac12u_0 .\nonumber
\end{align}
Here $e=\|E\|$ and $Q=\|K^{1/2}E\|$.
\end{theorem}
\begin{proof}
On a common smooth core, the product rule gives
\[
 i\langle v,[V,K]v\rangle
  =2\operatorname{Im}\sum_j
       \langle\sqrt{\lambda_j}D_jv,
                     \sqrt{\lambda_j}(\nabla_jV)v\rangle .
\]
Let $F=-i\sum_i(a_iD_i+\nabla_i a_i/2)$. In its written matrix order,
\[
 [D_j,F]=-i\sum_i
  \left((\nabla_ja_i)D_i+a_i[D_j,D_i]
                              +\tfrac12\nabla_j\nabla_i a_i\right).
\]
In the derivative of the kinetic form, the self-adjoint $F$ acting on
$D_jE$ cancels, leaving
\[
 -2\operatorname{Re}\sum_{ji}\lambda_j
       \langle D_jE,(\nabla_ja_i)D_iE\rangle
 -2\operatorname{Re}\sum_j\lambda_j\langle D_jE,A_jE\rangle .
\]
Their absolute values are at most $2LQ^2$ and
$2Q(G_0e+G_1Q)$. Also
$i[F,U]=\sum_i a_i\partial_iU$, whose quadratic form is bounded by
$Q(u_0e+u_1Q)$. Thus, retaining the forcing,
\[
 (Q^2)'\leq2\alpha Q^2+2\beta eQ
                         +2Q\|K^{1/2}r\|.
\]
Regularize $Q$ at zero and use Duhamel's norm bound for $e$.
The integrating factor proves \eqref{p4s:eq:force}.

For form data, set $R_\epsilon=(I+\epsilon K)^{-1}$ and
$K_\epsilon=K(I+\epsilon K)^{-1}$. The ordered identities
\[
 K_\epsilon'=R_\epsilon K'R_\epsilon,\qquad
 [V+F,K_\epsilon]=R_\epsilon[V+F,K]R_\epsilon
\]
give the same inequality with $Q_\epsilon=\|K_\epsilon^{1/2}E\|$.
Indeed $\|R_\epsilon E\|\leq e$,
$\|K^{1/2}R_\epsilon E\|\leq Q_\epsilon$, and the forcing is bounded
by $Q_\epsilon\|K^{1/2}r\|$. Integrate before taking the monotone
spectral limit $\epsilon\downarrow0$. The assumed conforming
approximation supplies the propagator and commutator passage for
unbounded coefficients. Resolvent algebra alone does not prove those
domain premises.
\end{proof}

When $a_i=0$, this reduces to the spatial-force estimate with
$\alpha=a_K/2+g_1$ and $\beta=g_0$. A covariantly constant carrier
$M(t)$ has $\nabla_iM=0$ and contributes no spatial force. Its phase
is still in the true evolution. A laboratory-constant matrix can instead
have $[C_i,M]\ne0$ and must then be priced. For a scalar dilation
$a=cx$ in one flat coordinate, the kinetic contribution is
$-2c\|\partial_xE\|^2$, confirming the commutator sign. Spatially
varying mobility or a discontinuous coefficient needs additional
terms or a different theorem.

\subsection{A positive invariant when a quadratic trap is inverted}

\begin{proposition}[Transported quadratic form and physical coflow]
\label{p4s:prop:invariant}
Let $H_0(t)$ be a real scalar Weyl-ordered quadratic rate Hamiltonian
in $z=(q,p)$, $p=-i\partial_q$, including its affine terms. Let $S(t)$
be its classical symplectic propagator and $m(t)$ its affine classical
solution. For any positive matrix $W_0$, set
\[
 W(t)=S(t)W_0S(t)^T,\qquad
 K(t)=1+\tfrac14(z-m)^TW(t)^{-1}(z-m).
\]
Then $K\geq1$ and, as forms on the transported Schwartz core,
\begin{equation}\label{p4s:eq:invariant}
 K_t+i[H_0,K]=0 .
\end{equation}
For a pure Gaussian with phase gradient
$m_p+R(q-m_q)$, $R=R^T$, and positive imaginary phase matrix $I$,
its covariance has blocks
\[
 W=\begin{pmatrix}C&CR\\RC&RCR+I/2\end{pmatrix},\qquad C=(2I)^{-1}.
\]
The invariant is therefore
\[
 K=1+\tfrac12\left[
 (p-m_p-R(q-m_q))^TI^{-1}(p-m_p-R(q-m_q))
                     +(q-m_q)^TI(q-m_q)\right].
\]
Writing $D_0=\partial-i(m_p+R(q-m_q))$ and
$\Lambda=\tfrac12I^{-1}$, one has $K\geq D_0^\dagger\Lambda D_0$.
The Gaussian ground value of $K$ is $1+d/2$.
\end{proposition}
\begin{proof}
If the classical quadratic generator is $A=JG_0$, then
$W'=AW+WA^T$ and
$(W^{-1})'=-A^TW^{-1}-W^{-1}A$. The affine equation for $m$
cancels the linear terms. Commutators of Weyl quadratics equal their
Poisson expressions, since higher Moyal derivatives vanish; these
identities prove \eqref{p4s:eq:invariant}. A positive quadratic
matrix is a sum of squares of real linear canonical observables,
proving positivity. Factoring the displayed covariance as
\[
 \begin{pmatrix}1&0\\R&1\end{pmatrix}
 \begin{pmatrix}C&0\\0&I/2\end{pmatrix}
 \begin{pmatrix}1&R\\0&1\end{pmatrix}
\]
and inverting proves the form formula. Each centered Gaussian
coordinate contributes $1/2$ to its quadratic ground value.
\end{proof}

This is established invariant mathematics
\cite{Lewis1967,LewisRiesenfeld1969}; its role here is a positive
error metric through an inverted interval. It is dimensionless and
does not replace the physical mass in the current. For the same
laboratory vector potential $A_0$, define
$b_{0,i}=(\hbar m_{p,i}+\hbar(R(q-m_q))_i-e_iA_{0,i})/M_i$.
For any true wave in that constitution,
\[
 j_i[\Psi]-\rho_\Psi b_{0,i}
       =\frac{\hbar}{M_i}\operatorname{Im}(\Psi^\dagger D_{0,i}\Psi).
\]
Consequently its difference from the analogous expression for $P$ is
at most
\begin{equation}\label{p4s:eq:invariantcurrent}
 \frac{\hbar}{M_i}
 \left(e\|D_{0,i}P\|
             +N\sqrt{(\Lambda^{-1})_{ii}}\,Q\right)
\end{equation}
in integrated absolute value. The weighted coordinate Cauchy
inequality proves the last factor. A changed potential $\Delta A$
adds $-e_i\Psi^\dagger\Delta A_i\Psi/M_i$ to the relative current.

For $H=H_0+V$ with Hermitian multiplication $V$ and form forcing $r$,
the exact cancellation \eqref{p4s:eq:invariant} gives
$Q'\leq g_1Q+g_0e+\|K^{1/2}r\|$ if
$\|\Lambda^{1/2}\nabla V\,v\|\leq g_0\|v\|+g_1\|K^{1/2}v\|$.
The proof is the preceding commutator calculation, with no separate
$K_t$ loss. Resolvent regularization preserves
$K_{\epsilon,t}+i[H_0,K_\epsilon]=0$.
A canonical drift rewritten in $D_0$ also creates the multiplication
term $a\cdot(m_p+R(q-m_q))$; its force and literal current remain.
An internal-branch-dependent covariance would produce further
commutators and is outside this scalar invariant statement.

\subsection{Temporal transfer with a retained internal clock}

\begin{theorem}[First-domain causal estimate with a form-controlled time derivative]
\label{p4s:thm:temporal}
Let $H(t)$ be self-adjoint on a fixed domain $\mathcal D$, with
$t\mapsto H(t)\in\mathcal L(\mathcal D,\mathcal H)$ continuously
differentiable for an equivalent fixed graph norm. Assume its unitary
common-domain propagator is locally bounded in that graph norm.
Let $D_s$ be a fixed finite Hermitian right-clock matrix, and write
$\mathcal L(t)X=H(t)X-XD_s$.
For an absolutely continuous Hilbert--Schmidt residual $r$, with
$r'\in L^1$, consider
\[
 i\delta'=\mathcal L\delta-r,\qquad \delta(0)=0.
\]
Suppose $K=H+b(t)\geq I$ and, on $\mathcal D$,
\[
 \|(H'-\omega')v\|\leq c_1(t)\|K^{1/2}v\|+c_0(t)\|v\|
\]
for a declared real differentiable scalar $\omega$ and integrable
nonnegative $c_0,c_1$. For each $T$ define
\begin{align*}
 E_T&=\int_0^T\|r\|,\quad E_{D,T}=\int_0^T\|rD_s\|,
       \quad B_T=\sup_{[0,T]}|b+\omega|,\\
 C_T&=B_TE_T^2+E_TE_{D,T},\qquad G_T=\int_0^Tc_1,\\
 A_T&=\sup_{0\leq t\leq T}
 \left[\|r(t)\|+\|r(0)\|+\int_0^t
       \left(\|r'+i\omega r\|+c_0(s)\int_0^s\|r(u)\|\,du\right)ds\right].
\end{align*}
When these quantities are finite, the causal solution belongs to the
common strong domain and satisfies
\begin{equation}\label{p4s:eq:temporal}
 \sup_{[0,T]}\|K^{1/2}\delta\|
       \leq\sqrt{E_TA_T+C_T}+\tfrac12E_TG_T .
\end{equation}
No spatial form norm of $r$ and no application of $H$ twice to
$\delta$ is assumed.
\end{theorem}
\begin{proof}
We first justify the domain, rather than use a formal $H^2$ calculation.
Set $B(t)=(\mathcal L(t)-i)^{-1}$. It maps the Hilbert--Schmidt space
boundedly into its common graph domain and
$B'=-B\mathcal L'B$. If $U(t,s)$ is the two-sided propagator, integration
by parts in the graph norm gives
\begin{align}
 \delta(t)={}&B(t)r(t)-U(t,0)B(0)r(0)\nonumber\\
 &+\int_0^tU(t,s)
       \left(Br+B\mathcal L'Br-Br'\right)(s)\,ds .
       \label{p4s:eq:graphconstruction}
\end{align}
Indeed $\partial_sU(t,s)=iU(t,s)\mathcal L(s)$ and
$i\mathcal L B=i-B$, so differentiating $UBr$ recovers
$iU r$ minus the displayed integrand. This is precisely the
Duhamel error. The terms are graph-integrable; absolutely continuous
approximation in time proves the identity for the stated data.

Put $\eta=(\mathcal L-\omega)\delta-r$. Testing against
common-domain adjoint solutions, or using
\eqref{p4s:eq:graphconstruction}, yields the mild identity
\[
 i\eta'=\mathcal L\eta-i(r'+i\omega r)
                           +i(H'-\omega')\delta,\qquad\eta(0)=-r(0).
\]
It does not assert $\eta\in\mathcal D$. Let
$\ell=\|(\mathcal L-\omega)\delta\|$, $e=\|\delta\|$ and
$Q=\|K^{1/2}\delta\|$. Unitarity and the commuting right multiplication
give
$e(t)\leq E(t)=\int_0^t\|r\|$ and
$\|\delta D_s\|\leq E_D(t)=\int_0^t\|rD_s\|$.
The auxiliary equation and
$K\delta=(\mathcal L-\omega)\delta+(b+\omega)\delta+\delta D_s$
therefore imply
\[
 \ell(t)\leq A_T+\int_0^tc_1Q,\qquad
 Q(t)^2\leq E_T\ell(t)+C_T .
\]
For $E_T>0$, set $M(t)=A_T+\int_0^tc_1Q$. Then
$M'\leq c_1\sqrt{E_TM+C_T}$, and differentiation of the square root,
with a positive regularizer if necessary, gives
$\sqrt{E_TM+C_T}\leq\sqrt{E_TA_T+C_T}+E_T\int_0^tc_1/2$.
This proves \eqref{p4s:eq:temporal}. If $E_T=0$, the causal error is
identically zero.
\end{proof}

The common-domain evolution premise is a standard nonautonomous
evolution requirement; the regularity equivalence in
\cite{SchmidGriesemer2014} is relevant to checking it. The estimate
retains the actual right-clock defect $E_{D,T}$ and the covariant
temporal combination $r'+i\omega r$. A large internal operator is
not a removable scalar phase.

For example, with
$\Pi_j=p_j-e_j(A_{0,j}+a_j(t))$ in physical units and rate generator
$H_{\rm kin}=\sum_j\Pi_j^2/(2M_j\hbar)$,
\[
 H_{\rm kin}'=
 -\sum_j\frac{e_j}{M_j\hbar}a_{j,t}\cdot\Pi_j
 +i\sum_j\frac{e_j}{2M_j}\operatorname{div}a_{j,t}.
\]
If $K$ dominates the complete kinetic form, bounded vector additions
give the admissible coefficients
\[
 c_1^2=\sum_j\frac{2e_j^2\|a_{j,t}\|_\infty^2}{M_j\hbar},
 \quad
 c_0=\sum_j\frac{|e_j|\|\operatorname{div}a_{j,t}\|_\infty}{2M_j}
       +\|(V_{\rm add}/\hbar)'-\omega'I\|_\infty .
\]
A linear force $F_j(t)x_j$ is also allowed when fixed positive traps
provide a position form: Young's inequality supplies a shift
$b\geq1+\|V_{\rm add}^-\|/\hbar+
\sum_jF_j^2/(M_j\Omega_j^2\hbar)$ and domination of
$M_j\Omega_j^2x_j^2/(4\hbar)$.
Its additional contribution to $c_1^2$ is
$\sum_j4|F_j'|^2/(M_j\Omega_j^2\hbar)$.
A changing quadratic trap or an unbounded source-coordinate field
needs its actual stronger graph bound.

\subsection{Singular residuals and a safe common-domain alternative}

\begin{proposition}[Coulomb residual and temporal coordinates]
\label{p4s:prop:coulomb}
In three relative spatial dimensions, let $\phi$ be a smooth Gaussian
nonzero at a collision. Then $|x|^{-1}\phi\in L^2_{\rm loc}$ but
$|x|^{-1}\phi\notin H^1_{\rm loc}$. Nevertheless, if $V(x)=g/|x|$
is fixed in laboratory coordinates, $P_2(t,x)$ is a differentiable
quadratic comparison potential and $\phi_t$ is a polynomial times a
smooth finite Gaussian family, then
\[
 f=(V-P_2)\phi,\qquad
 f_t=-P_{2,t}\phi+(V-P_2)\phi_t
\]
are locally in $L^2$ and globally in $L^2$ for the indicated Gaussian
polynomial families. On a finite parameter path with nondegenerate
covariance they are continuous in $L^2$.
Under a homogeneous point change $x=L(t)X$, with invertible $L$,
\[
 |\partial_t(g/|L(t)X|)|
       \leq\|L'L^{-1}\|\,|g|/|L(t)X|.
\]
\end{proposition}
\begin{proof}
The radial square integrals near zero have orders
$\int_0^1dr$ for $1/r$ and $\int_0^1r^{-2}dr$ for its leading
gradient. The nonzero Gaussian value prevents cancellation of the
latter leading term. Polynomial factors and Gaussian tails give the
global $L^2$ statements; local domination by an integrable $r^{-2}$
square envelope and uniform Gaussian tail bounds prove continuity.
Finally differentiation gives
$-g\,x\cdot(L'L^{-1}x)/|x|^3$, proving the last inequality.
\end{proof}

An exponentially small nonzero Gaussian value does not remove the
divergent spatial derivative. A translation that moves the collision
point instead produces a $1/r^2$ temporal multiplier. The homogeneous
repair retains the physical centroid, linear force, boost and scalar
action; recentering them away may restore that moving singularity.
These observations select a valid response method, without claiming a
small apparatus error.

For completeness a strong-graph variant applies when a changing trap
has a graph-relative rather than form-relative time derivative.
\begin{proposition}[One graph by an ordered inverse lift]
\label{p4s:prop:graph}
Let $H(t)$ satisfy the common-domain assumptions above, and let
$A=H+c\geq aI$ for constants $a>0,c\geq0$.
Assume $\gamma(t)=\|H'(t)A(t)^{-1}\|\in L^1$.
For $iE'=HE+f$, $E(0)\in\mathcal D$ and
$f\in W^{1,1}([0,T];\mathcal H)$, set $R(t)=\int_0^t\gamma$.
Then
\begin{align}
 \|A(t)E(t)\|\leq{}&\|f(t)\|+
 e^{R(t)}\biggl[\|A(0)E(0)\|+\|f(0)\|\nonumber\\
 &\hspace{9mm}+\int_0^t e^{-R(s)}
       \bigl((c+\gamma(s))\|f(s)\|+\|f'(s)\|\bigr)\,ds\biggr].
       \label{p4s:eq:graphlift}
\end{align}
\end{proposition}
\begin{proof}
The auxiliary $Y=AE+f$ obeys
$iY'=HY+cf+iH'E+if'$; the terms containing $Hf$ cancel.
Solve its Volterra equation using
$E=A^{-1}(Y-f)$ and the bounded perturbation $H'A^{-1}$.
The norm inequality for $Y$ and Gronwall give
\eqref{p4s:eq:graphlift}. This is also a rigorous construction:
$(A^{-1})'=-A^{-1}H'A^{-1}$ verifies the original equation weakly
for the constructed $E$, and uniqueness of the Hilbert mild solution
identifies it. No separate expression $Hf$ is needed.
\end{proof}

For a fixed laboratory Coulomb singularity and moving oscillator
centres, the time derivative of the Coulomb kernel is zero. The
derivative of the trap is quadratic and linear in the laboratory
coordinates, so it can be graph-relative under the earned oscillator
domain. Proposition~\ref{p4s:prop:graph}, rather than an unjustified
spatial derivative of $f$, is then the relevant interface. A detector
moving with those centres still pays every term in
\eqref{p4s:eq:surface}.

\section{A complete Coulomb forcing measure}
\label{p4h:sec:spectral}

This application keeps both the electron relative position and a
quantum source coordinate. Its first calculation concerns a fixed
Coulomb reference, whose exact spectral measure can be evaluated.
The second restores the actual source interaction and feedback.
These are different estimates: a continuum-only coefficient of the
fixed reference cannot remove the poles or replace the coupled
resolvent of the actual parent.

Use reduced-mass atomic units and write
\begin{equation}\label{p4h:eq:electron}
 h=-\tfrac12\Delta_r-\frac1{|r|},\qquad
 \phi(r)=\pi^{-1/2}e^{-|r|},\qquad f(r)=r_z\phi(r).
\end{equation}
The Coulomb operator is self-adjoint on $H^2(\mathbb R^3)$ with
form domain $H^1$, ground energy $-1/2$ and normalized ground
$\phi$. Its excited bound energies are $E_n=-1/(2n^2)$,
$n\geq2$, with continuum $[0,\infty)$.
We use the usual Coulomb spectrum, not a newly inferred set of levels.
Only the isolated forcing $f$ has the exact angular selection
$l=1,m=0$; the actual coupled parent below retains every angular sector.

\begin{theorem}[Full spectral measure of the dipole forcing]
\label{p4h:thm:spectral}
The spectral measure of $f$ under $h$ is
\begin{equation}\label{p4h:eq:measure}
 \mu_f=\sum_{n=2}^\infty d_n^2\delta_{-1/(2n^2)}
                       +1_{E>0}\rho_C(E)\,dE,
\end{equation}
where
\begin{align}
 d_n^2&=\frac{256n^7(n-1)^{2n-5}}{3(n+1)^{2n+5}},
 \label{p4h:eq:atoms}\\
 \rho_C(k^2/2)&=\frac{256}{3}
 \frac{\exp[-4\arctan(k)/k]}
 {(1+k^2)^5(1-e^{-2\pi/k})},\qquad k>0 .
 \label{p4h:eq:continuumdensity}
\end{align}
Its total mass and first $(h+1)$ form moment are both one:
\begin{equation}\label{p4h:eq:sumrules}
 \int d\mu_f=1,\qquad \int E\,d\mu_f=0,\qquad
                    \int(E+1)\,d\mu_f=1 .
\end{equation}
On the whole true continuum,
\begin{equation}\label{p4h:eq:densityceiling}
 \rho_C(E)\leq\rho_*,\qquad
 (1+E)\rho_C(E)\leq\rho_*,
 \qquad \rho_*=\frac{256}{3e^4}.
\end{equation}
\end{theorem}
\begin{proof}
With $\psi=u(r)Y_{10}(\widehat r)/r$, the radial forcing and
Friedrichs operator are
\[
 u_f=\frac2{\sqrt3}r^2e^{-r},\qquad
 h_1=-\tfrac12\partial_r^2+r^{-2}-r^{-1}.
\]
Integration of $r^4e^{-2r}$ gives $\|f\|^2=1$.
Direct differentiation gives
$h_1u_f=(-1/2+1/r)u_f$, whence $\langle f,hf\rangle=0$.

We specify the continuum normalization because an energy
cross-section convention could otherwise introduce a wrong Jacobian.
The regular Coulomb function is normalized by its unit-amplitude
large-radius oscillation. In the standard convention
\cite{DLMFCoulomb},
\begin{align*}
 u_k(r)&=\sqrt{2/\pi}\,F_1(-1/k,kr),\\
 F_1(-1/k,kr)&=C_1(kr)^2e^{-ikr}\,
                    {}_1F_1(2+i/k;4;2ikr),\\
 C_1^2&=\frac{2\pi(1+k^2)}
                  {9k^3(1-e^{-2\pi/k})}.
\end{align*}
For $H_1^+=G_1+iF_1$, the Wronskian in $r$ is
$W_r(F_1,H_1^+)=-k$. The upper-boundary radial resolvent kernel is
\begin{equation}\label{p4h:eq:Green}
 (h_1-E-i0)^{-1}(r,r')
   =\frac2k F_1(-1/k,kr_<)H_1^+(-1/k,kr_>).
\end{equation}
Its derivative jump is $-2$, as required by the coefficient
$-1/2$ in $h_1$. Its imaginary part is $(2/k)F_1(r)F_1(r')$.
Stone's formula therefore gives the measure $u_ku_k\,dk$,
or $u_ku_k/k$ per unit energy $E=k^2/2$.
The Coulomb asymptotics used to choose $H_1^+$ are asymptotic
statements at large radius, not exact finite-radius equalities.

Here is the overlap calculation. Put
$a=2+i/k$, $s=1+ik$, $z=2ik/(1+ik)$.
Laplace integration of the confluent hypergeometric series,
first in its absolute-convergence region and then by analytic
continuation, gives
\[
 \int_0^\infty r^4e^{-sr}{}_1F_1(a;4;2ikr)\,dr
       =24s^{-5}{}_2F_1(a,5;4;z).
\]
The coefficient identity $(5)_j/(4)_j=1+j/4$ implies
\[
 {}_2F_1(a,5;4;z)
 =(1-z)^{-a}\left(1+\frac{az}{4(1-z)}\right).
\]
The last bracket is $1/[2(1-ik)]$. Consequently the integral is
\[
 \frac{12}{(1+ik)^5(1-ik)}
       \left(\frac{1-ik}{1+ik}\right)^{-2-i/k}.
\]
The continuous logarithm of the ratio is $-2i\arctan k$.
Its contribution to the squared modulus is
$e^{-4\arctan(k)/k}$. Multiplication by the normalized radial
prefactors gives
$|\langle u_k,u_f\rangle|^2=k\rho_C(k^2/2)$.
The factor $1/k$ from Stone's energy measure proves
\eqref{p4h:eq:continuumdensity}.

The normalized negative-energy eigenfunction in this channel is
\[
 u_{n1}(r)=\frac{4r^2}{n^3}
 \sqrt{\frac{(n-2)!}{(n+1)!}}\,
        e^{-r/n}L_{n-2}^3(2r/n).
\]
Insert $L_j^3=(4)_j\,{}_1F_1(-j;4;\cdot)/j!$ into the same
Laplace identity. The bracket is now $n/[2(n-1)]$ and gives
\eqref{p4h:eq:atoms}; in particular $d_2^2=32768/59049$.

For completeness, the radial regular/decaying solutions have
only the simple negative poles $-1/(2n^2)$, $n\geq2$.
Their residues give precisely these normalized eigenprojections.
On every compact positive-energy interval,
\eqref{p4h:eq:Green} has continuous boundary values against the
exponentially decaying forcing, so its forcing measure there
is absolutely continuous. A remaining singular measure could
only be supported at zero. The regular zero-energy solution is
proportional to $\sqrt r\,J_3(\sqrt{8r})$; its large-radius
oscillatory magnitude is of order $r^{1/4}$ and it is not $L^2$.
Thus zero is not an eigenvalue and supplies no atom.
A singular continuous measure cannot be supported on a single
point. This proves completeness of \eqref{p4h:eq:measure}.
The already evaluated norm and first energy expectation now
give \eqref{p4h:eq:sumrules}; numerical quadrature is not used
to establish them.

To prove the whole-continuum ceiling, let
$A(k)=\arctan k-k/(1+k^2)$. Since
\[
 \frac{d}{dk}\left\{\frac{2k^3}{3(1+k^2)}-A(k)\right\}
             =\frac{2k^4}{3(1+k^2)^2}\geq0,
\]
we have $A(k)\leq2k^3/[3(1+k^2)]$.
The logarithmic derivative of $(1+k^2/2)\rho_C(k^2/2)$
is consequently at most
\[
 -\frac{16k}{3(1+k^2)}
       +\frac{2\pi}{k^2(e^{2\pi/k}-1)}.
\]
For $k\leq1$, use $e^x-1\geq x^3/6$ to bound the positive
term by $3k/(2\pi^2)$; this leaves a strictly negative margin.
For $k\geq1$, use $e^x-1>x$ to bound it by $1/k$,
whereas the negative magnitude is at least $8/(3k)$.
The threshold limit is $\rho_*$. This proves the weighted
ceiling and hence the unweighted one.
\end{proof}

The bound lines are substantial. The formula gives
$n^3d_n^2\longrightarrow\rho_*$ and
$E_{n+1}-E_n\sim n^{-3}$, matching the threshold density.
For $n\geq100$ one also has
\[
 d_n^2\leq\frac7{4n^3},\qquad
              \sum_{n>N}d_n^2\leq\frac7{8N^2}.
\]
For example the logarithm of the factor multiplying $256/(3n^3)$
is at most $-4+10/n\leq-3.9$, and
$256e^{-3.9}/3<7/4$ by a positive exponential Taylor sum.
These tails can check a truncated calculation; they do not
turn the discrete spectrum into a continuum.

\subsection{A continuum-only coherent response}

Let
\begin{equation}\label{p4h:eq:source}
 H_s=\frac{\Omega}{2}(-\partial_Y^2+Y^2),\qquad
 U_s(t)=e^{-itH_s},\qquad\Omega>0 .
\end{equation}
The source input is a vector in its physical-coordinate Hilbert
space, possibly tensored with an inert finite internal reference.

\begin{proposition}[Fixed-reference continuum response]
\label{p4h:prop:continuum}
For the fixed parent $h+H_s$ and forcing $f\otimes a(t)$,
let $q_C(0)=0$ denote the electronic continuum response. Then
\begin{align}
 \|q_C(t)\|^2,\quad\|(h+1)^{1/2}q_C(t)\|^2
 &\leq2\pi\rho_*\int_0^t\|a(s)\|^2\,ds,
 \label{p4h:eq:continuumresponse}\\
 \|H_s^{1/2}q_C(t)\|^2
 &\leq2\pi\rho_*\int_0^t\|H_s^{1/2}a(s)\|^2\,ds
 \label{p4h:eq:continuumsource}
\end{align}
whenever the displayed inputs are finite.
\end{proposition}
\begin{proof}
In the cyclic spectral representation generated by $f$, the
continuum wave equals
\[
 q_C(E,t)=-iF(E)e^{-iEt}U_s(t)
       \int_0^t e^{iEs}U_s(-s)a(s)\,ds,\qquad
                         |F(E)|^2=\rho_C(E).
\]
Extend the time input by zero outside $[0,t]$. Hilbert-valued
Plancherel, integrated first on the whole frequency line,
gives the factor $2\pi$ for this unnormalized Fourier integral.
Restricting to $E>0$ and using
\eqref{p4h:eq:densityceiling} gives both estimates in
\eqref{p4h:eq:continuumresponse}. The source form commutes its
own propagation, so the same argument applied to
$H_s^{1/2}a$ gives \eqref{p4h:eq:continuumsource}.
\end{proof}

For $a=gYp$, the last input is $gH_s^{1/2}Yp$, in that order.
It is not $gYH_s^{1/2}p$. The source interaction picture here
is used only for the norm proof; it has not changed the
physical source coordinate or its current. The bound poles
are excluded from this proposition and must be retained
separately: a Poisson-smoothed atom contributes
$d_n^2/(\pi\eta)$ at its center. There is no uniform
$\eta$-independent full-spectrum density ceiling.

\section{The actual compact-profile electron--source parent}
\label{p4h:sec:actual}

Choose a real radial cutoff with
$\chi_R=1$ for $|r|\leq R$, $\chi_R=0$ for $|r|\geq R+1$,
$0\leq\chi_R\leq1$ and $|\chi_R'|\leq8$. Such a smooth
flat cutoff is explicit: on $0<t<1$ use
\[
 \chi_R(R+t)=1-
 \frac{e^{-1/t}}{e^{-1/t}+e^{-1/(1-t)}}.
\]
For $t\leq1/2$ its derivative is bounded by
$e^{2-1/t}(t^{-2}+4)\leq8$; reflection treats the other half.
Set
\begin{equation}\label{p4h:eq:parent}
 d_R(r)=r_z\chi_R(|r|),\quad d_0=\|d_R\|_\infty\leq R+1,\qquad
 H=h+H_s+gYd_R ,
\end{equation}
with real constant $g$. This is a prescribed compact-profile
scalar parent. The finite annular electrostatic density
$-\epsilon_0\Delta d_R$ does not by itself construct its
quantized source, field energy or laboratory realization.

\begin{lemma}[Physical domains and forcing tail]
\label{p4h:lem:parent}
The parent \eqref{p4h:eq:parent} is self-adjoint on
$D(h)\cap D(H_s)$ and has form domain
$H^1(\mathbb R^3\times\mathbb R)\cap\{Y\psi\in L^2\}$.
It satisfies
\begin{equation}\label{p4h:eq:parentfloor}
 H+2+\frac{g^2d_0^2}{\Omega}
             \geq I+\tfrac14p_r^2+\tfrac12H_s .
\end{equation}
For $f_R=d_R\phi$ and $R\geq40$,
\begin{align}
 \|f_R-f\|^2&\leq\epsilon_R^2:=
 e^{-2R}\left(\tfrac23R^4+\tfrac43R^3+2R^2+2R+1\right),
 \label{p4h:eq:tailnorm}\\
 \|(h+1)^{1/2}(f_R-f)\|
 &\leq\sqrt{41.5+R^{-2}}\,\epsilon_R .
 \label{p4h:eq:tailform}
\end{align}
At $R=40$ these upper bounds are respectively
$6\cdot10^{-15}$ in norm and $4\cdot10^{-14}$ in first-form norm.
\end{lemma}
\begin{proof}
Hardy and interpolation make Coulomb infinitesimally
Laplacian bounded. The independent electron and oscillator
operators commute; after a common positive shift their sum
has domain $D(h)\cap D(H_s)$.
Since
$\|Y\psi\|\leq\sqrt{2/\Omega}\|H_s^{1/2}\psi\|$ and $d_R$
is bounded, $gYd_R$ is infinitesimally operator bounded
relative to this sum. This proves the operator assertion and
the corresponding common closed form.

The square
$\frac14\|\nabla u+2\widehat r\,u\|^2\geq0$ gives
$h+1\geq p_r^2/4$. Young's inequality gives
$gYd_R\geq-\Omega Y^2/4-g^2d_0^2/\Omega$ and
$H_s-\Omega Y^2/4\geq H_s/2$, proving
\eqref{p4h:eq:parentfloor}.

The omitted forcing is confined to $r\geq R$ and has magnitude
at most that of $u_f$. Integrating
$(4/3)r^4e^{-2r}$ gives \eqref{p4h:eq:tailnorm}.
For the radial error $w=(\chi_R-1)u_f$,
$|u_f'|\leq u_f$ on this tail and therefore $|w'|\leq9u_f$.
The positive upper bound on the remaining radial form
potential is $1+r^{-2}$. Thus
\[
 \langle w,(h_1+1)w\rangle
       \leq\left(\tfrac{81}2+1+R^{-2}\right)
                          \int_R^\infty|u_f|^2\,dr,
\]
which is \eqref{p4h:eq:tailform}. Reciprocal positive Taylor
sums for $e^{80}$ enclose the stated numerical tails.
\end{proof}

The cutoff is essential to the actual source parent.
With an uncut $gYr_z$, completing the source square leaves
$-g^2r_z^2/(2\Omega)$. Electron packets translated to
$r_z=L$ and source packets translated to $Y=-gL/\Omega$
have energy tending to $-\infty$. Thus no stable uncut
source Hamiltonian is being approximated here.
Equations \eqref{p4h:eq:tailnorm}--\eqref{p4h:eq:tailform}
compare forcing vectors, not complete actual parents.

\subsection{Exact blocks, including source feedback}

Let
\[
 P=|\phi\rangle\langle\phi|\otimes I,\qquad Q=I-P,\qquad
 \Psi=\phi p+q,\quad q=Q\Psi .
\]
The projection $P$ does not reduce $H$.
The odd profile has $Pd_RP=0$, so the exact equations are
\begin{align}
 i\dot q&=Aq+Lp,& A&=QHQ,& Lp&=gf_R\otimes Yp,
 \label{p4h:eq:qblock}\\
 i\dot p&=(H_s-\tfrac12)p+
                        gY\langle\phi,d_Rq\rangle_r .
 \label{p4h:eq:pfeedback}
\end{align}
The subscript on the last inner product means integration in
the electron coordinate only. The actual vector $p(Y,t)$
retains its source phase and every finite internal coefficient.
It is not prescribed as a free source waveform.

On $Q$ put
\begin{equation}\label{p4h:eq:weights}
 A_0=Q(h+H_s)Q,\qquad K_0=A_0+1,\qquad K=A+1,\qquad
 T_{\rm ph}=I+\tfrac12p_r^2+H_s .
\end{equation}
Here $A_0$ is the self-adjoint restriction of $h+H_s$ to
$Q\mathcal H$. The spectral projection $Q$ preserves the common
operator and form domains. More precisely, $A$ denotes the
self-adjoint operator $A_0+gYQd_RQ$ on $D(A_0)$: the electronic
factor $Qd_RQ$ is bounded and commutes with the source coordinate,
so the same infinitesimal relative-bound argument used in
Lemma~\ref{p4h:lem:parent} applies. Thus $QHQ$ denotes this proved
realization, rather than an assumed self-adjoint compression.

Since the complete excited electronic floor is $-1/8$,
\begin{equation}\label{p4h:eq:K0bounds}
 K_0\geq7/8+H_s,\qquad K_0\geq p_r^2/4+H_s,\qquad
                         T_{\rm ph}\leq(22/7)K_0 .
\end{equation}
For the last inequality take $4/11$ of the first lower
bound and $7/11$ of the second, then multiply by $22/7$.
This gives the exact coefficients $1$ and $1/2$ of the
identity and electron kinetic terms; the source coefficient
is larger than required.

The unbounded source interaction has the bounded form sandwich
\begin{equation}\label{p4h:eq:theta}
 \|K_0^{-1/2}gYQd_RQK_0^{-1/2}\|
 \leq\theta:=|g|d_0\sqrt{\frac{2}{\Omega(7/8)}}.
\end{equation}
Indeed use $\|u\|\leq(7/8)^{-1/2}\|K_0^{1/2}u\|$ on
one factor, and
$\|Yv\|\leq\sqrt{2/\Omega}\|K_0^{1/2}v\|$ on the other.
This bounds the full sesquilinear form and hence its operator.
If $\theta<1$, then
\begin{equation}\label{p4h:eq:Kcomparison}
 (1-\theta)K_0\leq K\leq(1+\theta)K_0 .
\end{equation}
All these inequalities are physical form inequalities on the
complete $Q$ space. They neither truncate the oscillator nor
delete $QYd_RQ$.

\subsection{The spectral supremum and its correction}
\label{p4h:sec:gap}

For $\zeta=E+i\eta$, $\eta>0$, the spectral theorem gives
\begin{align}
 J_{\rm spec}(E,\eta)
 &:=\|K_0^{1/2}(A_0-\zeta)^{-1}K_0^{1/2}\|
   =\sup_{a\in\sigma(A_0)}\frac{a+1}{|a-\zeta|}
 \notag\\
 &\leq\overline J(E,\eta):=
     \sup_{a\geq-1/8+\Omega/2}\frac{a+1}{|a-\zeta|}.
 \label{p4h:eq:Jspectrum}
\end{align}
The inequality can be strict. The half-line in the last
expression contains spectral gaps and cannot generally
replace $\sigma(A_0)$ in an equality.

\begin{example}[A strict gap in the half-line bound]
\label{p4h:ex:gap}
Take $\Omega=.01$, $E=-.117$, $\eta=.0001$.
The lowest spectral points of $A_0$ are $-.120$ and $-.110$;
the $n\geq3$ electronic ladders begin above $-.051$, and
the continuum begins at $.005$. At $a=E$, the half-line
expression is $8830$. On the actual spectrum the maximum is
\[
 J_{\rm spec}=\frac{.88}{\sqrt{.003^2+.0001^2}}<294.
\]
Indeed the first point gives a value above $280$, while for
every other spectral point $a\geq-.110$ the quotient is
at most $.89/.007<128$. Thus the half-line equality fails
by a large factor even for this concrete parent.
\end{example}

\begin{theorem}[An earned low-window inverse for the actual block]
\label{p4h:thm:lowwindow}
If $\overline J\theta<1$, then
\begin{equation}\label{p4h:eq:lowgeneral}
 \|K_0^{1/2}(A-\zeta)^{-1}LH_s^{-1/2}\|
 \leq\frac{|g|}{1-\overline J\theta}
 \left\{\frac2\Omega\sup_{s\geq\Omega/2}
 \int\frac{e+s+1}{|e+s-\zeta|^2}\,\mu_{f_R}(de)\right\}^{1/2}.
\end{equation}
For the static row
\begin{equation}\label{p4h:eq:parameters}
 \Omega=1/100,\quad |g|\leq10^{-6},\quad R=40,\quad d_0\leq41,
 \qquad E\leq-.47,\quad\eta>0,
\end{equation}
the complete physical sandwich satisfies
\begin{equation}\label{p4h:eq:lowrow}
 \|T_{\rm ph}^{1/2}(A-E-i\eta)^{-1}LH_s^{-1/2}\|^2
                  <4.035635\cdot10^{-9}<4.05\cdot10^{-9}.
\end{equation}
\end{theorem}
\begin{proof}
Let $V_0=K_0^{-1/2}gYQd_RQK_0^{-1/2}$ and
$B_0=K_0^{1/2}(A_0-\zeta)^{-1}K_0^{1/2}$.
The bounded form pencil of $A-\zeta$ has inverse
\[
 K_0^{1/2}(A-\zeta)^{-1}K_0^{1/2}
           =(I+B_0V_0)^{-1}B_0.
\]
The order is fixed: the pencil is
$B_0^{-1}+V_0=B_0^{-1}(I+B_0V_0)$.
The Neumann bound is
$\|(I+B_0V_0)^{-1}\|\leq(1-\overline J\theta)^{-1}$.
This is the concrete form-inverse construction of
Theorem~\ref{p4f:thm:ordered}; no Hilbert-norm boundedness of
$YQd_RQ$ has been assumed.

In the baseline squared norm, integrate out only the electron
spectral measure. The remaining positive source operator is
\[
 H_s^{-1/2}YF(H_s)YH_s^{-1/2},\qquad
 F(s)=\int\frac{e+s+1}{|e+s-\zeta|^2}\,\mu_{f_R}(de).
\]
Bound $F$ by its supremum, then use
$\|YH_s^{-1/2}\|\leq\sqrt{2/\Omega}$.
This proves \eqref{p4h:eq:lowgeneral}.
Moving $Y$ through $F(H_s)$ or through the resolvent would
not prove that estimate.

For \eqref{p4h:eq:parameters},
$\theta<.00062$ and $\overline J\leq88/35$.
To see the latter, first lower every denominator to
$a+.47$ and use $a\geq-.12$; the resulting
$(a+1)/(a+.47)$ decreases and has value $88/35$ at the floor.
Thus the correction in \eqref{p4h:eq:Jspectrum} leaves the
low-window bound unchanged.

The source supremum also needs its own argument.
First dominate the integral kernel at the fixed endpoint
$E=-.47,\eta=0$, since $e+s\geq-.12$.
At this endpoint the derivative of
$(e+s+1)/(e+s+.47)^2$ with respect to $s$ has sign
$-(e+s+1.53)<0$. It is consequently enough to use $s=1/200$.
This monotonicity was not asserted for every negative $E$
before making the endpoint domination.

For the uncut forcing, isolate the exact $2p$ mass
$w_2=32768/59049$. All remaining spectral mass lies at
$e\geq-1/18$, and the same endpoint kernel is decreasing
in $e$. Its complete bound-plus-continuum integral is therefore
\begin{equation}\label{p4h:eq:Fceiling}
 F_f(1/200)\leq
 w_2\frac{352}{49}+(1-w_2)\frac{123048}{22801}<6.4 .
\end{equation}
This analytic ceiling requires neither a continuum quadrature
nor a truncation of the Rydberg series.

The compact forcing changes the baseline weighted inverse
norm, before multiplication by $|g|$, by at most
\[
 \overline J\,\epsilon_R\sqrt{\frac2{\Omega(7/8)}}.
\]
This follows by inserting $K_0^{-1/2}$ and applying its
lower form floor to $f_R-f$ and the ordered source factor.
Use \eqref{p4h:eq:K0bounds}, $\epsilon_R<6\cdot10^{-15}$,
$\sqrt{1280}<35.778$ and $\sqrt{1600/7}<16$.
The resulting entirely rational upper bound is
\[
 \frac{22}{7}\,10^{-12}
 \frac{\big[35.778+(88/35)(6\cdot10^{-15})16\big]^2}
      {[1-(88/35)(.00062)]^2}
       <4.035635\cdot10^{-9}.
\]
This proves \eqref{p4h:eq:lowrow}, retaining the actual
excited-block feedback in the inverse.
\end{proof}

The result is a uniform frequency-window statement with source
weight $H_s^{1/2}$. It does not say that the actual time-dependent
input $p$ has Fourier support in this window. Such a use requires
the temporal band and complementary-frequency terms of the
window theorem. We next give a different all-frequency bound
with the stronger, explicitly earned source weight $H_s$.

\section{Actual source graphs and finite-time complete currents}
\label{p4h:sec:currents}

Write $m_0=\|f_R\|^2$ and
$m_1=\|(h+1)^{1/2}f_R\|^2$. The form tail gives
$m_0\leq1$ and, at $R=40$, $m_1<(1+4\cdot10^{-14})^2
<1+10^{-12}$.
The source ladder operators give the ordered bounds
\begin{equation}\label{p4h:eq:ladder}
 \|YH_s^{-1}\|\leq\frac{4\sqrt2}{3\Omega},\qquad
 \|H_s^{1/2}YH_s^{-1}\|
       \leq\frac{\sqrt3+1/\sqrt2}{\sqrt\Omega}.
\end{equation}
Indeed $H_s=\Omega(N+1/2)$ and $Y=(a+a^\dagger)/\sqrt2$.
The downward and upward coefficients of $aH_s^{-1}$ and
$a^\dagger H_s^{-1}$ are bounded by
$2/(3\Omega)$ and $2/\Omega$, respectively.
For the weighted versions, their squared ratios are
\[
 \frac{n(n-1/2)}{\Omega(n+1/2)^2}\leq\frac1\Omega,\qquad
 \frac{(n+1)(n+3/2)}{\Omega(n+1/2)^2}\leq\frac6\Omega.
\]
The first inequality is direct and the second ratio decreases
from its value $6$ at $n=0$. Their sum in $Y$ proves
\eqref{p4h:eq:ladder}. These are full-source operator bounds,
not a bound on a selected actual occupation.

Consequently
\begin{equation}\label{p4h:eq:Comega}
 \|K_0^{1/2}(f_R\otimes Y)H_s^{-1}\|^2
 \leq C_\Omega^2:=
 \frac{32m_1}{9\Omega^2}
 +\frac{(\sqrt3+1/\sqrt2)^2m_0}{\Omega}.
\end{equation}
The two terms follow by splitting
$K_0=(h+1)+H_s$ on the forcing vector, with the ordering
in \eqref{p4h:eq:ladder} preserved.

\begin{proposition}[All-frequency response with a physical source graph]
\label{p4h:prop:global}
For $\theta<1$ and every $E\in\mathbb R$, $\eta>0$,
\begin{equation}\label{p4h:eq:globalinverse}
 \|T_{\rm ph}^{1/2}(A-E-i\eta)^{-1}LH_s^{-1}\|^2
 \leq \frac{22}{7}\frac{1+\theta}{1-\theta}
                         \frac{g^2C_\Omega^2}{\eta^2}.
\end{equation}
For the actual input $p$ of \eqref{p4h:eq:pfeedback}, with
$q(0)=0$, let $q_T$ agree with its forced $Q$ equation on
$[0,T]$ and then evolve homogeneously with $A$. If
$H_sp\in L^2(0,T)$, then
\begin{equation}\label{p4h:eq:damped}
 \int_0^\infty e^{-2\eta t}\|T_{\rm ph}^{1/2}q_T(t)\|^2dt
 \leq \frac{22}{7}\frac{1+\theta}{1-\theta}
       \frac{g^2C_\Omega^2}{\eta^2}
                \int_0^T e^{-2\eta s}\|H_sp(s)\|^2ds .
\end{equation}
Also, whenever $H_sp\in L^1(0,t)$,
\begin{equation}\label{p4h:eq:stationary}
 Q(t):=\|T_{\rm ph}^{1/2}q(t)\|
 \leq |g|\sqrt{\frac{22}{7}\frac{1+\theta}{1-\theta}}\,
                         C_\Omega\int_0^t\|H_sp(s)\|\,ds .
\end{equation}
\end{proposition}
\begin{proof}
By \eqref{p4h:eq:K0bounds} and \eqref{p4h:eq:Kcomparison},
$T_{\rm ph}\leq22K/[7(1-\theta)]$.
The form $K=A+1$ commutes with the actual $A$ resolvent,
whose unweighted norm is at most $1/\eta$.
Its input satisfies
$\|K^{1/2}LH_s^{-1}\|
\leq\sqrt{1+\theta}|g|C_\Omega$.
Their ordered product proves \eqref{p4h:eq:globalinverse}.
The finite lower bound on $K$ is positive; rescaling by that
bound gives the equivalent unit-coercive form convention
if needed. No commutation of $T_{\rm ph}$ with $A$ is used.

Extend $1_{[0,T]}H_sp$ by zero. The causal exponentially
damped Fourier transform of the Duhamel equation is the
resolvent multiplier applied to this complete source vector.
Plancherel with \eqref{p4h:eq:globalinverse} proves
\eqref{p4h:eq:damped}, as in Theorem~\ref{p4f:thm:damped}.
There is no derivative of the time cutoff in this argument
and no frequency-band assumption.
Finally Duhamel in the stationary positive $K$ form and
$\|K^{1/2}e^{-itA}v\|=\|K^{1/2}v\|$ give
\eqref{p4h:eq:stationary}.
\end{proof}

The $1/\eta^2$ cost includes all actual bound poles and
continuum. It cannot be discarded because the fixed-reference
continuum-only calculation happened to be $\eta$ independent.
The source weights in \eqref{p4h:eq:lowgeneral} and
\eqref{p4h:eq:globalinverse} are also different.

\begin{proposition}[Source graphs for the actual feedback evolution]
\label{p4h:prop:actualsource}
Let the normalized entrance be
$\Psi(0)=\phi\otimes|1\rangle\otimes\zeta$, with any normalized
passive finite internal vector $\zeta$. Under the actual static
parent \eqref{p4h:eq:parent}, put $a=|g|d_0$.
Then
\begin{align}
 H_1(t):=\|H_s^{1/2}\Psi(t)\|
 &\leq\sqrt{3\Omega/2}+a\sqrt{\Omega/2}\,t,
 \label{p4h:eq:Hone}\\
 \|H_s\Psi(t)\|
 &\leq\Omega\{3/2+\sqrt3\,at+a^2t^2/2\},
 \label{p4h:eq:Htwo}\\
 e(t):=\|q(t)\|
 &\leq |g|\{\sqrt3\,t+at^2/2\}.
 \label{p4h:eq:enorm}
\end{align}
Since $P$ commutes with $H_s$, the first two bounds also
control the corresponding graphs of the actual $p$.
\end{proposition}
\begin{proof}
The commutator is
$[H_s,H]=-ig\Omega d_Rp_Y$, with $p_Y=-i\partial_Y$.
The expectation identity and
$\|p_Y\Psi\|\leq\sqrt{2/\Omega}\|H_s^{1/2}\Psi\|$ give
\[
 \left|\frac d{dt}H_1^2\right|
      \leq a\sqrt{2\Omega}\,H_1 .
\]
Regularizing the square root if necessary and integrating
gives \eqref{p4h:eq:Hone}.
The norm derivative for $H_s\Psi$ obeys
\[
 \frac d{dt}\|H_s\Psi\|
   \leq\|[H_s,H]\Psi\|
   \leq a\sqrt{2\Omega}\,H_1(t).
\]
Its initial norm is $3\Omega/2$; integrating the preceding
linear bound gives \eqref{p4h:eq:Htwo}.

These computations can be justified without assuming the
future source stays Gaussian. Let $P_N$ project onto source
occupations $0,\ldots,N-1$, and use the full-space operators
\[
 H_N=h+H_s+g d_RP_NYP_N.
\]
They are self-adjoint on $D(h)\cap D(H_s)$ with uniform
infinitesimal relative bounds and the coercivity
\eqref{p4h:eq:parentfloor}. For every vector in this common domain,
$H_Nu\to Hu$, because $P_N$ commutes with $H_s$ and converges
in its form norm. The resolvent identity, with the uniformly bounded
nonreal resolvents on its left, gives strong resolvent convergence;
the associated unitary groups converge strongly on compact time
intervals. For $N\geq2$ the entrance lies in $P_N\mathcal H$ and
remains there, so both source graphs are bounded operators on that
invariant subspace. The commutator is
$-ig\Omega d_RP_Np_YP_N$ and obeys the same uniform estimates as
above. Hence the displayed graph bounds hold for $e^{-itH_N}\Psi(0)$.
At each time, strong wave convergence and weak lower semicontinuity
of the closed graphs of $H_s^{1/2}$ and $H_s$ pass these bounds to
the actual wave. No finite source truncation remains in the result.

Lastly, from the exact equation for $q$ and its zero
entrance,
\[
 \|q(t)\|\leq |g|\|f_R\|\int_0^t\|Yp(s)\|\,ds
 \leq |g|\sqrt{2/\Omega}\int_0^tH_1(s)\,ds .
\]
Here $\|f_R\|\leq1$ and $P$ commutes with $H_s$.
This is \eqref{p4h:eq:enorm}. The feedback term in
\eqref{p4h:eq:pfeedback} has not been set to zero anywhere.
\end{proof}

\subsection{All coherent terms in the original physical currents}

The nominated current constitution is the nonrelativistic
canonical one:
\[
 j_r[\Psi]=\operatorname{Im}\Psi^\dagger\nabla_r\Psi,\qquad
 j_Y[\Psi]=\Omega\operatorname{Im}\Psi^\dagger\partial_Y\Psi.
\]
Define differences relative to the \emph{actual} $P$ wave,
\[
 \delta\rho=|\Psi|^2-|\phi p|^2,\qquad
 \delta j=j[\Psi]-j[\phi p].
\]
The vector $\phi p$ is generally unnormalized and does not
solve an autonomous closed evolution. It is a wave
comparison, not a separately assigned physical flow. To make its
continuity defect explicit, set
$\mathcal F=gY\langle\phi,d_Rq\rangle_r$ and $M(t)=\|p(t)\|^2$.
The feedback equation gives, in distributions,
\begin{equation}\label{p4h:eq:comparisondefect}
 \partial_t|\phi p|^2+\operatorname{div}j[\phi p]
 =s_P:=2|\phi|^2\operatorname{Im}(p^\dagger\mathcal F),
 \qquad M=1-\|q\|^2,\quad M'=\int s_P.
\end{equation}
The products are integrable under the earned source graphs: for
example $\|\mathcal F\|\leq |g|d_0\|Yq\|$.
Even if $M>0$ and the comparison is normalized, writing
$\rho_P=|\phi p|^2$, its density $\rho_P/M$ and current $j[\phi p]/M$ have source
$s_P/M-M'\rho_P/M^2$. Normalization therefore does not create
conservative equivariance. The estimates below retain the original
unnormalized comparison and need neither a second flow nor a
continuity-defect allowance.

\begin{lemma}[Full electron and source current differences]
\label{p4h:lem:current}
With $d_r=\|\nabla_rq\|$ and $d_Y=\|\partial_Yq\|$,
\begin{align}
 \|\delta\rho\|_1&\leq2e+e^2,\label{p4h:eq:densitydifference}\\
 \|\delta j_r\|_1&\leq d_r+e+ed_r,\label{p4h:eq:jrelectron}\\
 \|\delta j_Y\|_1
 &\leq\Omega\{(1+e)d_Y+
                     e\sqrt{2/\Omega}\,H_1(t)\},
 \label{p4h:eq:jYsource}
\end{align}
where
$d_r\leq\sqrt2\,Q(t)$ and
$d_Y\leq\sqrt{2/\Omega}\,Q(t)$.
\end{lemma}
\begin{proof}
For any derivative $D$, expand before estimating:
\[
 j_D[\phi p+q]-j_D[\phi p]
 =c_D\operatorname{Im}\{
     (\phi p)^\dagger Dq+q^\dagger D(\phi p)+q^\dagger Dq\}.
\]
Thus every interference term and the quadratic error current
remain. Cauchy--Schwarz, $\|p\|\leq1$ and
$\|\nabla\phi\|=1$ give the electron estimate.
For the source,
$\|\partial_Yp\|\leq\sqrt{2/\Omega}\,H_1$ because $P$
commutes with the source kinetic form. Multiplication by
$c_Y=\Omega$ gives \eqref{p4h:eq:jYsource}.
The density expansion gives \eqref{p4h:eq:densitydifference}.
Finally the electron and source kinetic terms are included
with their stated coefficients in $T_{\rm ph}$, giving the
two derivative bounds. Spectral orthogonality of $P$ and $Q$
has not been used to discard a pointwise cross term.
\end{proof}

\begin{corollary}[An explicit finite-time current bound]
\label{p4h:cor:currents}
For the same static parent with
$\Omega=.01$, $g=10^{-6}$, $R=40$, $d_0\leq41$ and
the entrance of Proposition~\ref{p4h:prop:actualsource},
\begin{align}
 Q(100)&<.00051,\qquad e(100)<.000176,\label{p4h:eq:qrow}\\
 \int_0^{100}\|\delta j_r(t)\|_1dt
 &<.045033131<.046,\label{p4h:eq:electronrow}\\
 \int_0^{100}\|\delta j_Y(t)\|_1dt
 &<.003970830<.004.\label{p4h:eq:sourcerow}
\end{align}
\end{corollary}
\begin{proof}
Equations \eqref{p4h:eq:tailform}, \eqref{p4h:eq:theta}
and \eqref{p4h:eq:Comega} give the rational enclosure
\[
 \frac{22}{7}\frac{1+.00062}{1-.00062}
 \left\{\frac{32(1+10^{-12})}{9(.01)^2}
                       +\frac{25}{4(.01)}\right\}<338^2 .
\]
Here $(\sqrt3+1/\sqrt2)^2=7/2+\sqrt6<25/4$.
Let $a=41\cdot10^{-6}$. Equations
\eqref{p4h:eq:stationary}--\eqref{p4h:eq:enorm} consequently
admit the polynomial upper bounds
\begin{align*}
 \overline Q(t)&=338\cdot10^{-8}
           \{(3/2)t+(7/8)at^2+a^2t^3/6\},\\
 \overline e(t)&=10^{-6}\{(7/4)t+at^2/2\},\\
 \overline H_1(t)&=.123+.071at.
\end{align*}
All roundings are outward:
$\sqrt3<7/4$, $\sqrt{.015}<.123$ and
$\sqrt{.005}<.071$.
Substitute $d_r\leq(10/7)\overline Q$ and
$d_Y\leq15\overline Q$, using
$\sqrt2<10/7$, $\sqrt{200}<15$.
The two current integrands are bounded by
\[
 \frac{10}{7}\overline Q+\overline e+
                  \frac{10}{7}\overline e\,\overline Q,
 \qquad
 .01\{15(1+\overline e)\overline Q+
                                  15\overline e\,\overline H_1\}.
\]
These are explicit nonnegative rational polynomials.
Exact integration on $[0,100]$ gives respectively
$0.045033130836\ldots$ and $0.003970829749\ldots$.
Direct evaluation of the same polynomials gives
\eqref{p4h:eq:qrow}. The exact integrals are
\[
 \frac{113483489708747242649}{2520000000000000000000},
 \qquad
 \frac{95299913979542242649}{24000000000000000000000},
\]
so integer comparison verifies both printed outward ceilings.
No continuum quadrature enters this bound.
\end{proof}

The electron integral has units of relative-coordinate length,
and the source integral of source-coordinate length in the
chosen dimensionless $Y$ convention. They bound
$j[\Psi]-j[\phi p]$, not the entire absolute currents.
They are not event probabilities. A detector or tube claim
requires a corresponding surface/guard and an original-law
transport interface with its own constants. Comparing to a
freely propagated intended source instead of the actual $p$
would additionally require the feedback phase and current
cost of $p-p_{\rm free}$.

The compact scalar and oscillator constitute a definite
mathematical example with actual source feedback, full
Coulomb poles and continuum, and complete canonical currents.
A finite-mass source/COM reduction, source field energy,
retardation, loading, clock and a microscopic Pauli or Dirac
current identification remain separate physical questions.
None is inferred from the exact Coulomb calibration or from
the small forcing tail.

\subsection{Flow and companion-paper interfaces}
\label{p4h:sec:flowinterface}

For the entrance of Proposition~\ref{p4h:prop:actualsource}, the actual
parent \eqref{p4h:eq:parent} is the constant-$g$ case of the
\emph{Full Coulomb--source flow} theorem in the revised companion
paper~\cite{RodgersP3}. The reduced-mass units, the smooth compact
dipole and its full collar, the first excited source entrance, the
passive internal factor, and both canonical currents agree. Constant
$g$ satisfies that theorem's bounded absolutely continuous control
hypothesis on every finite horizon, including $[0,100]$.
Its common first-domain, tangent-regularity and logarithmic-current
arguments consequently supply a deterministic reference-compatible
flow for this actual wave, with reference-almost-sure collision and
node avoidance. An original joint law absolutely continuous with
respect to the entrance wave density is transported by reweighting
those same paths once. A quantitative domination factor for
\eqref{p4s:eq:path} is a further premise, not a consequence of absolute
continuity.

The current bounds above do not apply that flow theorem to $\phi p$.
Conservative two-flow stability in the revised P3 requires two
admitted conservative flows, a common positive tube, and a bound for
the probability of leaving it. Its separate signed-source appendix
requires a normalized reference, classical local regularity, and
finite source and divergence costs. Neither conclusion follows by
analogy from \eqref{p4h:eq:comparisondefect}. No variable-coefficient
spin-curl current is used here; adding one would require its own
derivative regularity and current estimates.

The Gaussian preparation and instrument results of P1~\cite{RodgersP1}
and the finite repeated-record protocol of P2~\cite{RodgersP2} concern
their own effective parents, complete entrance laws and decoders.
The present estimates establish neither those entrance premises nor
fresh independence, renewed readiness, or repeated unknown-input
measurement. Conversely, wave reset or preparation of an actual law
does not supply a missing derivative, source graph or flow hypothesis
for the present parent. Any composition must verify the same complete
configuration and every retained correlation on the same time interval.

\section{Significance and remaining scope}
\label{p4s:sec:conclusion}

The results give explicit routes from a complete coherent forcing to a
physical derivative form and then to every admitted coordinate current.
The inverse-lift argument handles a form-dual residual through its
temporal regularity. The window argument handles changing input
directions without discarding the operator response. The ordered
resolvent formulas retain source interactions and physical weights
without a false commutation. The Coulomb example discharges these
questions in a definite electron--oscillator parent, including the
source graph and the feedback-modified input.

The distinctions in the hypotheses are productive. They identify
whether a calculation still needs a residual derivative, a physical
form floor, a spectral coefficient, endpoint dressing, a source
moment, or a common-domain propagation argument. They also expose
specific invalid substitutions: a tiny collision tail does not make a
Coulomb residual differentiable in space; a source-only truncation
does not suppress propagated high-sector response; a reference
invariant is not the physical current metric; and a shifted spectral
floor is not an actual spectral gap.

A physical record application must still supply its original wave and
actual law, the complete Hamiltonian/current constitution, an earned
flow, a decoder and a whole retention interval. The current estimates
can then enter its guard or moving-surface argument with the relevant
length scales and original-law factor. No common apparatus joining
all four portfolio papers has been inferred. Finite material sources,
microscopic current reductions, preparation warrant and attainable
control errors remain separate research questions.

\paragraph{Funding.}
This research was conducted independently without external funding.

\paragraph{AI assistance and verification.}
AI systems assisted with mathematical drafting, literature comparison,
algebraic and numerical checks, and manuscript preparation. The author
remains responsible for the arguments and presentation. This assistance
does not constitute independent scientific validation.

\paragraph{Independent review and collaboration.}
The author welcomes independent mathematical review of the common-form
and domain arguments, computational replication of the explicit
spectral and current bounds, and physical collaboration on complete
source/control models and faithful retained records.

\appendix
\section{Spatial forcing compression with all sources retained}
\label{p4s:sec:spatial}

A bounded relative-position region can have discrete free spatial
modes even when the complete source Hilbert space has continuous
spectrum. The following construction truncates the spatial forcing
directions and leaves every source coefficient exact.

\begin{theorem}[Full-source inverse defect]
\label{p4s:thm:spatial}
Let $\mathcal H=L^2(\Omega;\mathcal Z)$, where $\mathcal Z$ includes
all source, centre-of-mass and internal variables. Let $a$ be a
closed positive form with operator $K$ and
\[
 a(v,v)\geq\|v\|^2+\lambda\|\nabla_{\rm rel}v\|^2
                 +\text{the required hidden kinetic forms}.
\]
Choose a bounded open relative-position domain $D\subset\Omega$, with
free Dirichlet eigenvalues $\mu_j$, and project onto its first $N$
spatial eigenfunctions tensored with $I_{\mathcal Z}$. In the full
physical form domain let $\mathcal W_N$ consist of zero extensions
of vectors in $H_0^1(D;\mathcal Z)$ whose coefficients in those $N$
modes vanish. This defines the zero Dirichlet condition without
requiring a classical boundary trace. Let
$\mathcal V_N=\mathcal W_N^{\perp_a}$ and let $K_N$ be the restricted
closed form operator on the Hilbert closure of $\mathcal V_N$.
With its physical inclusion $J_N$, define
\[
 B=K^{-1},\quad B_N=J_NK_N^{-1}J_N^\dagger,\quad D_N=B-B_N .
\]
Then $D_N\geq0$, its range is in $\mathcal W_N$, and
\begin{equation}\label{p4s:eq:spatial}
 \|D_N\|\leq\delta_N,\qquad
 \|K^{1/2}D_N\|\leq\sqrt{\delta_N},\qquad
 \delta_N=(1+\lambda\mu_{N+1})^{-1}.
\end{equation}
All exterior responses and Dirichlet form traces of $Bh$ and $B_Nh$
coincide. If $D\subset[-a,a]^3$ and $N=n^3$, a sufficient bound is
$\delta_N\leq[1+\lambda\pi^2(n+1)^2/(4a^2)]^{-1}$.
\end{theorem}
\begin{proof}
The coercive gradient bound makes form convergence imply convergence
in the relative $H^1$ norm. The zero extensions of
$H_0^1(D;\mathcal Z)$ form a closed subspace there; the spatial-mode
coefficient conditions are closed already in $L^2$. Thus
$\mathcal W_N$ and its form orthogonal complement are closed. Decompose the Riesz
solution $u=Bh$ into $u_N+w_N$. Testing its form equation in each
orthogonal subspace gives $u_N=B_Nh$, $w_N=D_Nh$ and
\[
 a(D_Nh,D_Nh)=\langle h,D_Nh\rangle .
\]
In particular $D_N$ is the positive inverse associated with the
restricted high form, included into the physical space. The vector
spectral expansion in the free spatial basis gives
$\|\nabla_{\rm rel}w\|^2\geq\mu_{N+1}\|w\|^2$ on $\mathcal W_N$.
Coercivity and Cauchy--Schwarz therefore yield
$\|D_Nh\|\leq\delta_N\|h\|$; substitution in the last identity
proves the form bound. Its support and zero Dirichlet trace give the
exterior claim; equality of conormal derivatives is not asserted.
Dirichlet domain monotonicity compares $D$ to the cube. Any cube
mode below $\pi^2(n+1)^2/(4a^2)$ has all three indices at most $n$,
so there are at most $n^3$ such modes, proving the stated ceiling.
\end{proof}

If the high form is dense in the corresponding high Hilbert space,
$D_N$ is its compressed-form inverse, included by zero extension.
Otherwise its Hilbert closure is used; the proof is unchanged.
It is generally \emph{not} the unrestricted inverse followed by a high
projection. For example, with
$K=\left(\begin{smallmatrix}2&1\\1&2\end{smallmatrix}\right)$ and
$\mathcal W=\operatorname{span}(0,1)$,
$D=\operatorname{diag}(0,1/2)$, whereas
$K^{-1}Q=\left(\begin{smallmatrix}0&-1/3\\0&2/3\end{smallmatrix}\right)$
is not even self-adjoint.

The decreasing high spaces have zero intersection by completeness of
the free spatial expansion. Their orthogonal projections in the form
Hilbert space converge strongly to zero: nested projection differences
have squared norms equal to the differences of the squared projection
norms, and the limit lies in the intersection. Hence the retained
spaces are form dense. One may enrich them by any finite list of
original form columns by intersecting the high space with their form
orthogonal complements. The defect decreases and
\eqref{p4s:eq:spatial} survives; an enriched $f\in\operatorname{Dom}K$
satisfies $B_NKf=f$. No actual population has been substituted.

The construction supplies an inverse error and exact source retention,
not a finite apparatus simulation: each spatial coefficient remains in
$\mathcal Z$, and boundary response spaces can be infinite dimensional.
Any evolution or control estimate must still propagate the physical
form under its own hypotheses. Once that estimate supplies $e,Q$,
Proposition~\ref{p4s:prop:current} controls every admitted coordinate
current, including source currents.

\section{A small omitted source does not remove causal feedback}
\label{p4s:sec:feedback}

\begin{proposition}[Retained high response and endpoint dressing]
\label{p4s:prop:feedback}
For an admitted same-parent block evolution, suppose
\[
 i u_L'=A_Lu_L+B^\dagger u_H+r_L,\qquad
 i u_H'=Du_H+Bu_L+r_H .
\]
Then, with the true high-block propagator $U_H$,
\[
 u_H(t)=U_H(t,0)u_H(0)
       -i\int_0^tU_H(t,s)\bigl(B(s)u_L(s)+r_H(s)\bigr)\,ds .
\]
Thus the low equation retains an entrance term, the high source,
and the retarded feedback kernel
$-iB(t)^\dagger\int_0^tU_H(t,s)B(s)u_L(s)\,ds$.
For a fixed $D\geq gI>0$ and absolutely continuous
$F=Bu_L+r_H$,
\begin{align}\label{p4s:eq:high}
 u_H(t)={}&e^{-itD}u_H(0)-D^{-1}F(t)
             +e^{-itD}D^{-1}F(0)\\
          &+D^{-1}\int_0^te^{-i(t-s)D}F'(s)\,ds .\nonumber
\end{align}
The last three terms have $D^{1/2}$ norm at most
$g^{-1/2}(\|F(t)\|+\|F(0)\|+\int_0^t\|F'\|)$.
\end{proposition}
\begin{proof}
The first formula is variation of constants in the true high block.
Substitution into the low equation gives the kernel and both other
terms. Since $\partial_se^{-i(t-s)D}=iDe^{-i(t-s)D}$,
integration by parts gives \eqref{p4s:eq:high}; the spectral theorem
bounds $\|D^{-1/2}\|\leq g^{-1/2}$.
\end{proof}

After removal of a constant real scalar frequency $\omega$, the relevant inverse
is $(D-\omega)^{-1}$ and the temporal quantity is $F'+i\omega F$.
A retained right clock must also be restored in that combination.
An angular kinetic gap in a frozen model does not by itself prove
the required gap of a full radial/source block.

The elementary two-level parent
$H=\left(\begin{smallmatrix}0&b\\b&\Delta\end{smallmatrix}\right)$
already isolates the issue. With original state $(1,0)$ and no forcing,
the high population is
\[
 \frac{4b^2}{\Delta^2+4b^2}
 \sin^2\!\left(\tfrac12t\sqrt{\Delta^2+4b^2}\right).
\]
This follows by subtracting $\Delta/2$ times the identity and squaring
the remaining traceless matrix. An exactly zero high source therefore
does not imply a zero high response. For a forced zero-entrance
solution with source $(g_0,0)$ one instead has
$u_H''(0)=-bg_0$, a separate causal counterexample.

In angular tensor applications a rank-two coupling connects both
top bands $l=L-1,L$ to omitted $L+1,L+2$ unless a genuine parity
restriction has been established. A concrete molecular application
must identify its full tensor coupling, correlated entrance and local
forcing bounds. Small forcing bounds do not replace
$Bu_L$, its propagated angular/source derivatives, or the original
entrance error in the formulas above. Coherent currents can respond
at first order in a high amplitude while its population is second
order: the current conversion must retain the cross terms even when
the high population is small.

\clearpage
\begin{thebibliography}{99}
\bibitem{RodgersP1}
J. Rodgers, \emph{Conditional Gaussian Preparation with Retained Archives},
companion manuscript (2026).
\href{https://doi.org/10.5281/zenodo.23259560}{doi:10.5281/zenodo.23259560}.

\bibitem{RodgersP2}
J. Rodgers, \emph{Effective Repeated Position Records with Retained Entropy
and Calibrated Reset}, companion manuscript (2026).
\href{https://doi.org/10.5281/zenodo.23259558}{doi:10.5281/zenodo.23259558}.

\bibitem{RodgersP3}
J. Rodgers, \emph{Reference-Weighted Deterministic Quantum Flows with
Singular Interactions and Retained Sources}, revised companion manuscript (2026).
\href{https://doi.org/10.5281/zenodo.23259569}{doi:10.5281/zenodo.23259569}.

\bibitem{RodgersEquilibrium}
J. Rodgers, \emph{Autonomous Quantum Measurement Chains with Faithful
Equilibrium Records}, consolidated revision of the version 2
massive-configuration preprint, 4 October 2026.
\href{https://doi.org/10.5281/zenodo.23131069}{doi:10.5281/zenodo.23131069}.

\bibitem{RodgersNonequilibrium}
J. Rodgers, \emph{Nonequilibrium Calibration and Faithful Records in
Autonomous Effective Measurement Models}, version 2, 4 October 2026 manuscript.
\href{https://doi.org/10.5281/zenodo.23131081}{doi:10.5281/zenodo.23131081}.

\bibitem{Kato1966}
T. Kato, Wave operators and similarity for some non-selfadjoint operators,
\emph{Mathematische Annalen} \textbf{162} (1966), 258--279.
\href{https://doi.org/10.1007/BF01360915}{doi:10.1007/BF01360915}.

\bibitem{DAncona2014}
P. D'Ancona, Kato smoothing and Strichartz estimates for wave equations with magnetic potentials,
\emph{Communications in Mathematical Physics} \textbf{335} (2015), 1--16.
\href{https://doi.org/10.1007/s00220-014-2169-8}{doi:10.1007/s00220-014-2169-8};
\href{https://arxiv.org/abs/1403.2537}{arXiv:1403.2537}.

\bibitem{Kisynski1964}
J. Kisy\'nski, Sur les op\'erateurs de Green des probl\`emes de Cauchy abstraits,
\emph{Studia Mathematica} \textbf{23} (1964), 285--328.
\href{https://doi.org/10.4064/sm-23-3-285-328}{doi:10.4064/sm-23-3-285-328}.

\bibitem{SchmidGriesemer2014}
J. Schmid and M. Griesemer, Kato's theorem on the integration of non-autonomous linear evolution equations,
\emph{Mathematical Physics, Analysis and Geometry} \textbf{17} (2014), 265--271.
\href{https://doi.org/10.1007/s11040-014-9154-5}{doi:10.1007/s11040-014-9154-5};
\href{https://arxiv.org/abs/1203.4700}{arXiv:1203.4700}.

\bibitem{DLMFCoulomb}
National Institute of Standards and Technology, \emph{Digital Library of Mathematical Functions},
\href{https://dlmf.nist.gov/33.2}{Section 33.2} and
\href{https://dlmf.nist.gov/33.11}{Section 33.11}, Coulomb functions.
Accessed 9 October 2026.

\bibitem{PanatiSpohnTeufel2003}
G. Panati, H. Spohn and S. Teufel, Space-adiabatic perturbation theory,
\emph{Advances in Theoretical and Mathematical Physics} \textbf{7} (2003), 145--204.
\href{https://doi.org/10.4310/ATMP.2003.v7.n1.a6}{doi:10.4310/ATMP.2003.v7.n1.a6};
\href{https://arxiv.org/abs/math-ph/0201055}{arXiv:math-ph/0201055}.

\bibitem{Lewis1967}
H. R. Lewis, Jr., Classical and quantum systems with time-dependent harmonic-oscillator-type Hamiltonians,
\emph{Physical Review Letters} \textbf{18} (1967), 510--512.
\href{https://doi.org/10.1103/PhysRevLett.18.510}{doi:10.1103/PhysRevLett.18.510};
erratum, \emph{Physical Review Letters} \textbf{18} (1967), 636,
\href{https://doi.org/10.1103/PhysRevLett.18.636.2}{doi:10.1103/PhysRevLett.18.636.2}.

\bibitem{LewisRiesenfeld1969}
H. R. Lewis, Jr. and W. B. Riesenfeld, An exact quantum theory of the time-dependent harmonic oscillator
and of a charged particle in a time-dependent electromagnetic field,
\emph{Journal of Mathematical Physics} \textbf{10} (1969), 1458--1473.
\href{https://doi.org/10.1063/1.1664991}{doi:10.1063/1.1664991}.
\end{thebibliography}

\end{document}
